\documentclass[11pt,a4paper]{book}
\usepackage[T1]{fontenc}
\usepackage[utf8]{inputenc}
\usepackage{lmodern}
\usepackage[english]{babel}
\usepackage{amsmath,amssymb,amsthm}
\usepackage{mathtools}
\usepackage{bbm}
\usepackage{booktabs}
\usepackage{longtable}
\usepackage{multirow}
\usepackage{array}
\usepackage{tabularx}
\usepackage{listings}
\usepackage{xcolor}
\usepackage{hyperref}
\usepackage[margin=1in]{geometry}
\usepackage{fancyhdr}
\usepackage{titlesec}
\usepackage{enumitem}
\usepackage{graphicx}
\usepackage{epigraph}
\setlength{\headheight}{14pt}

\hypersetup{
  colorlinks=true,
  linkcolor=blue,
  urlcolor=blue,
  citecolor=blue
}

\lstset{
  basicstyle=\ttfamily\small,
  breaklines=true,
  columns=fullflexible,
  frame=single,
  backgroundcolor=\color{gray!5}
}

\theoremstyle{definition}
\newtheorem{definition}{Definition}[chapter]
\newtheorem{theorem}{Theorem}[chapter]
\newtheorem{lemma}{Lemma}[chapter]
\newtheorem{corollary}{Corollary}[chapter]
\newtheorem{algorithm}{Algorithm}[chapter]
\newtheorem{example}{Example}[chapter]
\newtheorem{remark}{Remark}[chapter]

\pagestyle{fancy}
\fancyhf{}
\fancyhead[LE,RO]{\leftmark}
\fancyhead[RE,LO]{\rightmark}
\fancyfoot[C]{\thepage}

\titleformat{\chapter}[display]
  {\bfseries\Large}
  {\chaptername\ \thechapter}
  {20pt}
  {\Huge}

\setlist{nosep}

\begin{document}

\title{\Huge\textbf{Large Language Models: \\ A Multifaceted Mini-Book}\\[1cm]
	\large Information Theory, Deep Learning, Computer Theory,\\
	Software Engineering, Inference, and GPU Architecture}
\author{Bernard}
\date{\today}
\maketitle
\thispagestyle{empty}
\newpage

\tableofcontents
\newpage

% ============================================================
% Part I: Foundations
% ============================================================
\part{Foundations}

% ============================================================
\chapter{Information-Theoretic Foundations}
\label{ch:info-theory}

\epigraph{``The fundamental problem of communication is that of reproducing at one point either exactly or approximately a message selected at another point.''}{Claude Shannon, 1948}

This chapter develops the information-theoretic lens through which large language models (LLMs) are best understood. Language modelling is, at its core, a problem of efficient coding: a good language model assigns high probability to the next token, minimising the number of bits needed to encode the message.

\section{Entropy and Language}

\begin{definition}[Entropy]
	For a discrete random variable $X$ with distribution $p$, the Shannon entropy is
	\begin{equation}
		H(X) = -\sum_{x} p(x) \log_2 p(x) \quad\text{(bits)}.
	\end{equation}
	Entropy measures the \emph{minimum} number of bits required to encode $X$ on average.
\end{definition}

In natural language, the per-token entropy of English has been estimated at roughly $1.3$ bits/character (Shannon's game), though modern LLMs achieve cross-entropies well below $1$ bit/token on domain-specific corpora.

\begin{definition}[Cross-Entropy]
	Given a true distribution $p$ and a model $q$, the cross-entropy is
	\begin{equation}
		H(p,q) = -\sum_{x} p(x) \log_2 q(x) = H(p) + D_{\mathrm{KL}}(p \parallel q).
	\end{equation}
	Minimising cross-entropy is equivalent to minimising KL divergence from the true distribution.
\end{definition}

This is precisely the objective minimised during language model training. For a sequence $\mathbf{x} = (x_1, \dots, x_T)$, the per-token cross-entropy loss is
\begin{equation}
	\mathcal{L}_{\mathrm{NTP}} = -\frac{1}{T} \sum_{t=1}^{T} \log_2 p_\theta(x_t \mid x_{<t}),
\end{equation}
known as \emph{next-token prediction} (NTP).

\begin{theorem}[Shannon's Source Coding Theorem]
	For a stationary ergodic source with entropy rate $H$, the minimum expected code length per symbol achievable by any uniquely decodable code satisfies
	\begin{equation}
		H \leq L^* < H + 1.
	\end{equation}
\end{theorem}

This theorem provides the fundamental limit: no language model can achieve a cross-entropy lower than the true entropy of the language. In practice, LLMs approach but never reach this limit, and the gap is precisely the KL divergence.

\section{Kullback--Leibler Divergence}

\begin{definition}[KL Divergence]
	\begin{equation}
		D_{\mathrm{KL}}(p \parallel q) = \sum_{x} p(x) \log \frac{p(x)}{q(x)} \geq 0,
	\end{equation}
	with equality iff $p = q$ almost everywhere.
\end{definition}

KL divergence appears throughout LLM theory:
\begin{itemize}
	\item \textbf{Training:} minimising cross-entropy is minimising KL divergence.
	\item \textbf{Distillation:} student models are trained to minimise $D_{\mathrm{KL}}(p_\text{teacher} \parallel p_\text{student})$.
	\item \textbf{RLHF:} the KL penalty term $\beta D_{\mathrm{KL}}(\pi_{\mathrm{RL}} \parallel \pi_{\mathrm{SFT}})$ prevents the policy from diverging too far from the supervised model.
	\item \textbf{Beam search:} top-$k$ and top-$p$ sampling truncate the distribution, inducing controlled KL divergence.
\end{itemize}

\begin{theorem}[Gibbs' Inequality]
	For any two probability distributions $p, q$ over the same discrete space,
	\begin{equation}
		D_{\mathrm{KL}}(p \parallel q) \geq 0,
	\end{equation}
	with equality if and only if $p = q$ almost everywhere.
\end{theorem}

\begin{proof}
	By Jensen's inequality applied to the convex function $-\log$:
	\begin{equation}
		D_{\mathrm{KL}}(p \parallel q) = -\sum_x p(x) \log \frac{q(x)}{p(x)} \geq -\log \sum_x p(x) \frac{q(x)}{p(x)} = -\log 1 = 0.
	\end{equation}
\end{proof}

\section{Mutual Information in Attention}

\begin{definition}[Mutual Information]
	\begin{equation}
		I(X;Y) = D_{\mathrm{KL}}(p_{(X,Y)} \parallel p_X \otimes p_Y) = H(X) - H(X \mid Y).
	\end{equation}
	Mutual information measures the reduction in uncertainty about $X$ given $Y$.
\end{definition}

The attention mechanism computes a weighted sum of value vectors, where weights are given by a softmax over dot products:
\begin{equation}
	\alpha_{ij} = \frac{\exp(\mathbf{q}_i^\top \mathbf{k}_j / \sqrt{d_k})}{\sum_{j'} \exp(\mathbf{q}_i^\top \mathbf{k}_{j'} / \sqrt{d_k})}.
\end{equation}

This can be interpreted as a (normalised) pointwise mutual information between positions $i$ and $j$ in the latent space. The attention distribution $p(j \mid i)$ maximises the mutual information between the query position and the context positions under a soft constraint on the number of attended positions.

\begin{theorem}[Data Processing Inequality]
	If $X \to Y \to Z$ forms a Markov chain, then
	\begin{equation}
		I(X;Y) \geq I(X;Z).
	\end{equation}
\end{theorem}

In transformer layers, this implies that information about the input is monotonically non-increasing as it propagates through layers, though the \emph{task-relevant} information may increase through the information bottleneck.

\section{Rate--Distortion Theory and Quantisation}

Quantisation of LLM weights can be viewed through rate--distortion theory.

\begin{definition}[Rate--Distortion Function]
	For a source $X$ and distortion measure $d$, the rate--distortion function is
	\begin{equation}
		R(D) = \min_{p(\hat{x} \mid x): \mathbb{E}[d(X,\hat{X})] \leq D} I(X; \hat{X}).
	\end{equation}
	$R(D)$ gives the minimum number of bits per symbol needed to reconstruct the source with average distortion at most $D$.
\end{definition}

Post-training quantisation (PTQ) maps each weight $w \in \mathbb{R}$ to $\hat{w} \in \{0, \dots, 2^b-1\}$, typically with a scale and zero-point:
\begin{equation}
	\hat{w} = \text{clamp}\left(\left\lfloor \frac{w}{\Delta} \right\rceil + z, 0, 2^b-1\right).
\end{equation}
The distortion $d(w,\hat{w}) = (w - \hat{w})^2$ induces a degradation $D$ in model perplexity. Rate--distortion theory tells us that for a given bit-width $b$, there is a fundamental trade-off between compression rate and model quality.

\section{The Information Bottleneck and Hidden Representations}

The information bottleneck (IB) principle (Tishby et al., 1999) provides a framework for understanding transformer hidden states. An LLM layer computes a representation $Z$ that should be:
\begin{itemize}
	\item \emph{Sufficient:} $I(Z; Y) \geq I(X; Y)$ (preserve task-relevant information).
	\item \emph{Minimal:} minimise $I(Z; X)$ (compress away irrelevant details).
\end{itemize}

The IB Lagrangian is
\begin{equation}
	\mathcal{L}_{\mathrm{IB}} = I(Z; X) - \beta I(Z; Y).
\end{equation}

In transformers, each layer's hidden state $h_\ell$ can be seen as solving a successive IB trade-off, progressively filtering out task-irrelevant information while preserving predictive signal. This perspective explains why later layers are more task-specific and why pruning early layers often preserves performance.

\section{The Entropy of Natural Language}

Estimating the entropy of natural language is a longstanding problem:

\begin{itemize}
	\item \textbf{Shannon's game (1951):} humans guess the next letter; entropy $\approx 1.3$ bits/character.
	\item \textbf{Cross-entropy of n-gram models:} $\sim 1.75$ bits/char (5-gram on Brown corpus).
	\item \textbf{Neural language models:} GPT-2 achieves $\sim 1.0$ bits/char on WikiText-2.
	\item \textbf{LLMs:} GPT-4 achieves $\sim 0.7$ bits/char on carefully curated benchmarks.
\end{itemize}

The gap between model cross-entropy and the true entropy of the language ($H_\text{true} \approx 0.6$--$1.0$ bits/char) represents the irreducible loss due to fundamental unpredictability in natural language.

\section{Exercises}

\begin{enumerate}
	\item Prove that $H(p,q) \geq H(p)$ with equality iff $p = q$.
	\item Show that the cross-entropy loss for next-token prediction is an unbiased estimator of $H(p, p_\theta)$ in the limit of infinite data.
	\item For a quantised weight with $b$ bits, derive the expected MSE under uniform quantisation of a $\mathcal{N}(0, \sigma^2)$ distribution.
	\item Compute $R(D)$ for a Gaussian source with squared-error distortion.
	\item Explain why the KL penalty in RLHF prevents mode collapse.
	\item Prove the data processing inequality for Markov chains.
	\item Show that the IB Lagrangian can be rewritten as $\mathcal{L}_{\mathrm{IB}} = H(Y \mid Z) + \beta I(Z; X)$.
\end{enumerate}

% ============================================================
\chapter{Computational Theory of Language Models}
\label{ch:comp-theory}

\epigraph{``Computation is the process of transforming information according to well-defined rules.''}{Alan Turing}

This chapter examines LLMs through the lens of theoretical computer science: formal languages, computational complexity, and expressiveness.

\section{The Autoregressive Paradigm}

An autoregressive language model defines a distribution over sequences:
\begin{equation}
	p_\theta(x_1, \dots, x_T) = \prod_{t=1}^{T} p_\theta(x_t \mid x_{<t}).
\end{equation}

This is a \emph{generative} process: each token is sampled conditioned on all previous tokens. The computational cost is fundamentally sequential---generating token $t$ requires having generated all $t-1$ prior tokens.

\begin{theorem}[Sequential Lower Bound]
	Any exact autoregressive generation of a sequence of length $T$ requires at least $T$ sequential steps, regardless of parallel hardware.
\end{theorem}

\begin{proof}
	By induction. The distribution of $x_t$ depends on $x_{<t}$. Without knowledge of $x_{t-1}$, the marginal $p(x_t \mid x_{<t-1})$ differs from $p(x_t \mid x_{<t})$ unless the model is provably Markov of order $t-2$, which is not the case for general LLMs.
\end{proof}

This bound motivates non-autoregressive and speculative decoding techniques, which trade exactness for parallelism.

\section{Computational Complexity of Attention}

Let $n$ be the sequence length and $d$ the hidden dimension.

\begin{itemize}
	\item \textbf{Self-attention:} $O(n^2 d)$ time, $O(n^2)$ memory for the attention matrix.
	\item \textbf{Feed-forward:} $O(n d^2)$ time, $O(d^2)$ memory.
	\item \textbf{Autoregressive generation (per step):} $O(n d)$ for attention with KV cache (incremental), $O(d^2)$ for FFN.
\end{itemize}

\begin{theorem}[Context Window as Memory]
	The attention mechanism implements a content-addressable memory of size $O(n d)$. Retrieving information from a context of length $n$ requires $O(n)$ work, matching the lower bound for associative memory.
\end{theorem}

\begin{theorem}[Quadratic Lower Bound for Exact Attention]
	Computing the full attention matrix $A = \text{softmax}(Q K^\top / \sqrt{d_k})$ requires $\Omega(n^2)$ operations in the worst case, as every pair of positions $(i, j)$ must contribute to the output.
\end{theorem}

This quadratic barrier is the primary motivation for sparse attention patterns, linear attention, and SSM alternatives.

\section{Expressiveness: What Can Transformers Compute?}

\begin{theorem}[Universal Approximation for Sequence-to-Sequence]
	A transformer with sufficient width and depth can approximate any continuous sequence-to-sequence function on a compact domain, provided it has access to positional encodings.
\end{theorem}

\begin{theorem}[Turing Completeness]
	Transformers with rational-valued activations and CoT (chain-of-thought) prompting are Turing-complete. A transformer with $O(\log n)$ layers and polynomial width can simulate a Turing machine for $n$ steps.
\end{theorem}

\begin{corollary}[CoT Increases Expressiveness]
	There exist problems (e.g., evaluating Boolean formulas, solving linear equations) that cannot be solved by a constant-depth transformer without CoT but can be solved by the same transformer with $O(n)$ chain-of-thought steps.
\end{corollary}

\section{Circuit Complexity and TC0}

The class $\mathsf{TC}^0$ consists of problems solvable by constant-depth, polynomial-size circuits with threshold (majority) gates. Transformers with no chain-of-thought are believed to be limited to problems in $\mathsf{TC}^0$:

\begin{itemize}
	\item $\mathsf{TC}^0$ includes: parity, addition, multiplication by constants, sorting.
	\item $\mathsf{TC}^0$ does NOT include: multiplication of two $n$-bit numbers (requires $\mathsf{NC}^1$), matrix multiplication, or simulating a Turing machine for $n$ steps.
	\item With CoT, transformers can simulate $\mathsf{P}$ (polynomial time) by using the token sequence as a tape.
\end{itemize}

\begin{theorem}[Transformer Hierarchy]
	For a transformer with depth $L$, hidden dimension $d$, and precision $p$:
	\begin{itemize}
		\item Without CoT: computable functions lie in $\mathsf{TC}^0$ (for $L = O(1)$, $d = \text{poly}(n)$).
		\item With $O(\log n)$ CoT steps: $\mathsf{NC}^1$.
		\item With $\text{poly}(n)$ CoT steps: $\mathsf{P}$.
	\end{itemize}
\end{theorem}

\section{RASP: A Programming Language for Transformers}

RASP (Weiss et al., 2021) is a restricted programming language that captures the computational primitives of transformers:

\begin{itemize}
	\item \textbf{Selector operations:} attend from position $i$ to $j$ based on a predicate.
	\item \textbf{Aggregate:} compute a softmax-weighted average of values.
	\item \textbf{Sequence operations:} map, filter, reduce over positions.
\end{itemize}

Examples of RASP programs and their transformer realisations:
\begin{itemize}
	\item \textbf{Histogram:} count token frequencies in $O(1)$ layers.
	\item \textbf{Flip:} reverse a sequence in $O(1)$ layers.
	\item \textbf{dyck-1:} recognise balanced parentheses in $O(\log n)$ layers.
	\item \textbf{Prefix sum:} compute running sums in $O(1)$ layers.
\end{itemize}

\section{Injective LLMs and Input Reconstruction}

A recent line of work (arXiv:2510.15511) asks: can LLM-generated outputs be inverted to reconstruct the original input?

\begin{definition}[Injective Language Model]
	An LLM $p_\theta$ is injective if the mapping from input sequences to output distributions is one-to-one: different inputs always produce distinguishable output distributions.
\end{definition}

The SIPIT (Sampling-based Input Reconstruction) algorithm reconstructs the input prefix from generated tokens:

\begin{algorithm}[SIPIT]
	Given output tokens $y_1, \dots, y_m$, reconstruct input $x_1, \dots, x_n$ by:
	\begin{enumerate}
		\item Sampling multiple continuations from the model.
		\item Identifying input tokens that maximise the likelihood of observed outputs.
		\item Applying beam search over the input space.
	\end{enumerate}
\end{algorithm}

Empirically, SIPIT reconstructs up to 60\% of input tokens exactly for GPT-2, showing that LLMs encode substantial information about their inputs in their outputs.

\section{Formal Language Hierarchy}

LLMs correspond to different levels of the Chomsky hierarchy:

\begin{itemize}
	\item \textbf{Regular languages:} exactly expressible by bounded-length $n$-gram models (Markovian). A transformer without positional encodings is bounded to regular languages.
	\item \textbf{Context-free languages:} require stack-like memory; transformers with positional encodings can recognise some but not all CFLs. Dyck languages of depth $k$ require $\Omega(\log n)$ layers.
	\item \textbf{Context-sensitive languages:} require $O(n)$ memory; transformers with sufficient layers can recognise some CSLs via chain-of-thought.
	\item \textbf{Recursively enumerable:} Turing-complete with CoT.
\end{itemize}

\begin{theorem}[Separation by Depth]
	For any $k$, there exists a formal language $L_k$ recognisable by a transformer of depth $k+1$ but not by any transformer of depth $k$ (with bounded precision).
\end{theorem}

\section{The Scaling Hypothesis}

\begin{theorem}[Scaling Laws]
	For a transformer trained on $D$ tokens with $N$ parameters, the cross-entropy loss scales as
	\begin{equation}
		L(N, D) = E + \frac{A}{N^\alpha} + \frac{B}{D^\beta},
	\end{equation}
	where $E$ is the irreducible entropy of the data, $A,B,\alpha,\beta$ are constants, and $\alpha \approx \beta \approx 0.5$ for typical architectures.
\end{theorem}

The optimal allocation under a fixed compute budget $C$ satisfies $N_\text{opt} \propto C^{0.5}$ and $D_\text{opt} \propto C^{0.5}$ (Chinchilla scaling), implying model size and data should be scaled equally.

\begin{corollary}[Overtraining]
	Training a model beyond the Chinchilla-optimal point (more tokens than optimal) yields diminishing returns: doubling tokens gives at most $2^{-\beta} \approx 70\%$ of the loss reduction of optimally allocated compute.
\end{corollary}

\section{Embedding Condensation}

A recent finding (arXiv:2602.00217) shows that smaller models can ``condense'' the embedding geometry of larger models:

\begin{theorem}[Embedding Condensation]
	Training a smaller student model with a \emph{dispersion loss} $\mathcal{L}_{\text{disp}} = \|G_S - G_T\|_F^2$, where $G_S$ and $G_T$ are Gram matrices of hidden representations, transfers the geometric structure of the larger model's embedding space.
\end{theorem}

This suggests that the embedding geometry---the relative arrangement of token representations---carries substantial information that can be transferred across model scales.

\section{Exercises}

\begin{enumerate}
	\item Prove that exact autoregressive generation is in $\mathsf{P}$ but each step depends on all previous outputs.
	\item Show that the memory complexity of self-attention is $\Omega(n^2)$ for exact computation.
	\item Argue why a transformer without positional encodings cannot recognise the language $a^n b^n$.
	\item Derive the Chinchilla-optimal allocation from the scaling law equation.
	\item Explain why chain-of-thought can increase effective depth.
	\item Write a RASP program that computes the prefix sum of a binary sequence.
	\item Show that $\mathsf{TC}^0 \subseteq \mathsf{NC}^1$ and argue why transformer limitations follow.
\end{enumerate}

% ============================================================
% Part II: Architecture and Training
% ============================================================
\part{Architecture and Training}

% ============================================================
\chapter{Transformer Architecture}
\label{ch:transformer}

\epigraph{``Attention is all you need.''}{Vaswani et al., 2017}

\section{The Transformer Block}

Each transformer layer consists of two sub-layers:
\begin{enumerate}
	\item Multi-head self-attention (MHSA)
	\item Feed-forward network (FFN), typically a gated MLP
\end{enumerate}

Each sub-layer employs a residual connection and layer normalisation:
\begin{equation}
	h' = \text{LayerNorm}(h + \text{SubLayer}(h)).
\end{equation}

\subsection{Pre-Norm vs. Post-Norm}

The original transformer (Vaswani et al.) applied layer norm \emph{after} the residual addition (post-norm):
\begin{equation}
	h' = \text{LayerNorm}(h + \text{SubLayer}(h)).
\end{equation}

Modern LLMs (GPT-3, Llama, Mistral) apply layer norm \emph{before} each sub-layer (pre-norm):
\begin{equation}
	h' = h + \text{SubLayer}(\text{LayerNorm}(h)).
\end{equation}

Pre-norm provides more stable gradients at initialization, enabling training of deeper models without warmup instability.

\subsection{RMSNorm}

Llama and many subsequent models replace LayerNorm with RMSNorm:

\begin{align}
	\text{RMSNorm}(x)   & = \frac{x}{\sqrt{\frac{1}{d}\sum_{i=1}^d x_i^2}} \cdot \gamma,     \\
	\text{LayerNorm}(x) & = \frac{x - \mu}{\sqrt{\sigma^2 + \epsilon}} \cdot \gamma + \beta.
\end{align}

RMSNorm removes the mean-centering step, reducing computation by $O(d)$ per layer with negligible quality loss.

\section{Masked Multi-Head Attention}

In autoregressive decoding, each position can only attend to itself and earlier positions:
\begin{equation}
	\text{Attention}(Q, K, V) = \text{softmax}\left(\frac{QK^\top}{\sqrt{d_k}} + M\right)V,
\end{equation}
where $M$ is a causal mask: $M_{ij} = 0$ for $i \geq j$, $-\infty$ otherwise.

Multi-head attention splits the computation into $h$ heads:
\begin{equation}
	\text{head}_i = \text{Attention}(QW_i^Q, KW_i^K, VW_i^V),
\end{equation}
\begin{equation}
	\text{MHSA}(Q, K, V) = \text{Concat}(\text{head}_1, \dots, \text{head}_h) W^O.
\end{equation}

\section{Grouped Query Attention (GQA)}

Modern LLMs (Llama 2/3, Mistral) use GQA where multiple query heads share a key-value head:
\begin{equation}
	n_{\text{kv}} = \frac{n_{\text{q}}}{g}, \quad g \geq 1.
\end{equation}

GQA reduces KV cache memory by a factor of $g$ with minimal quality loss. It interpolates between multi-head attention ($g=1$) and multi-query attention ($g = n_q$).

\begin{theorem}[GQA Memory Savings]
	For a model with $L$ layers, $n_{kv}$ KV heads, head dimension $d_h$, sequence length $t$, and precision $p$ bytes, the KV cache memory is:
	\begin{equation}
		M_{\text{KV}} = 2 \cdot L \cdot n_{kv} \cdot t \cdot d_h \cdot p.
	\end{equation}
	With GQA ratio $g$, this becomes $M_{\text{KV}} / g$ compared to full MHA.
\end{theorem}

\section{Feed-Forward Networks}

The SwiGLU activation has become standard:
\begin{equation}
	\text{FFN}(x) = (W_1 x) \odot \sigma_{\text{SiLU}}(W_2 x) \cdot W_3,
\end{equation}
where $\sigma_{\text{SiLU}}(z) = z \cdot \sigma(z)$ is the SiLU/Swish activation, $\sigma$ is the logistic sigmoid, and $\odot$ denotes element-wise multiplication.

Other GLU variants:
\begin{itemize}
	\item \textbf{ReGLU:} $(W_1 x) \odot \max(0, W_2 x) \cdot W_3$.
	\item \textbf{GeGLU:} $(W_1 x) \odot \text{GELU}(W_2 x) \cdot W_3$.
	\item \textbf{SwiGLU} (default in Llama, Mistral, GPT-4).
\end{itemize}

For a hidden dimension $d$ and intermediate dimension $d_{ff} = \frac{8}{3}d$, SwiGLU has three weight matrices vs. two in the standard ReLU FFN, adding $33\%$ more parameters.

\section{RoPE: Rotary Positional Encodings}

RoPE encodes position via a rotation matrix:
\begin{equation}
	R_{\Theta,m} = \bigoplus_{j=1}^{d/2} \begin{pmatrix}
		\cos m\theta_j & -\sin m\theta_j \\
		\sin m\theta_j & \cos m\theta_j
	\end{pmatrix},
\end{equation}
where $\theta_j = 10000^{-2j/d}$.

The dot product $q_m^\top k_n$ in RoPE depends only on the relative position $m-n$:
\begin{equation}
	q_m^\top k_n = (R_{\Theta,m} x_q)^\top (R_{\Theta,n} x_k) = x_q^\top R_{\Theta, m-n} x_k.
\end{equation}

This relative-position property enables length generalisation. The rotating frequency $\theta_j$ determines how quickly position information decays: lower dimensions rotate faster (capturing local positions), higher dimensions rotate slower (capturing long-range positions).

\section{Attention Variants}

\begin{itemize}
	\item \textbf{ALiBi (Press et al., 2022):} adds a linear bias proportional to distance: $\text{softmax}(QK^\top / \sqrt{d} + m \cdot B)$ where $B_{ij} = -|i-j|$. No positional embeddings needed; extrapolates better to long sequences.
	\item \textbf{Attention with Linear Biases (xPos):} combines RoPE with exponential decay for better length generalisation.
	\item \textbf{Sparse attention:} only attend to local windows + global tokens (BigBird, Longformer).
\end{itemize}

\section{Exercises}

\begin{enumerate}
	\item Derive the FLOPs count per transformer layer.
	\item Show that GQA with $g$ groups reduces KV cache memory by a factor of $g$.
	\item Prove that RoPE makes attention depend only on relative position.
	\item Implement a single masked attention head in pseudocode.
	\item Compare pre-norm and post-norm in terms of gradient signal propagation.
	\item Compute the parameter count of a SwiGLU FFN vs. a standard ReLU FFN for the same $d_{ff}$.
\end{enumerate}

% ============================================================
\chapter{Tokenization}
\label{ch:tokenization}

\epigraph{``A tokenizer is a software stack, not a modelling choice.''}{Andrej Karpathy, 2023}

Tokenization converts raw text into a sequence of integers (token IDs) that the model processes. Despite being a preprocessing step, tokenization choices profoundly affect model behaviour, vocabulary efficiency, and multilingual capability.

\section{Byte-Pair Encoding (BPE)}

BPE (Gage, 1994; Sennrich et al., 2016) is the dominant tokenization algorithm, used by GPT-2/3/4, Llama, and most modern LLMs.

\begin{algorithm}[BPE Tokenizer Training]
	\begin{enumerate}
		\item Start with a vocabulary of individual bytes (or characters).
		\item Count all adjacent pairs of tokens in the corpus.
		\item Merge the most frequent pair into a new token.
		\item Repeat step 2--3 until the desired vocabulary size is reached.
	\end{enumerate}
\end{algorithm}

\begin{definition}[BPE Merge]
	A BPE merge operation replaces every occurrence of the adjacent pair $(a, b)$ in the corpus with the new symbol $ab$. The merge is \emph{greedy}: once merged, the pair becomes an atomic unit for all subsequent merges.
\end{definition}

The BPE tokenizer for GPT-2 operates on bytes (Byte-Level BPE), ensuring that any Unicode string can be tokenized without unknown tokens. Before BPE, most tokenizers operated on characters or word-level, requiring large vocabularies for multilingual coverage.

\begin{example}[BPE Merge Sequence]
	For the corpus $\{\text{``low''}, \text{``lower''}, \text{``lowest''}\}$:
	\begin{enumerate}
		\item Initial: \{\texttt{l}, \texttt{o}, \texttt{w}, \texttt{e}, \texttt{r}, \texttt{s}, \texttt{t}\}
		\item Merges: \texttt{l+o $\to$ lo}, \texttt{w+e $\to$ we}, \texttt{lo+w $\to$ low}, \texttt{we+r $\to$ wer}, \texttt{low+e $\to$ lowe}, \texttt{lowe+r $\to$ lower}, \texttt{low+e $\to$ lowe}, \texttt{low+s $\to$ lows}, \texttt{lowe+s $\to$ lowes}, \texttt{lowes+t $\to$ lowest}
	\end{enumerate}
	The final vocabulary size depends on the merge count limit.
\end{example}

\section{SentencePiece}

SentencePiece (Kudo \& Richardson, 2018) treats tokenization as a segmentation problem without needing pre-tokenized words (whitespace is handled as a regular symbol $\_$). It supports two algorithms:

\subsection{BPE Mode}

Same merge-based algorithm as BPE, but operating on the raw Unicode byte stream rather than pre-tokenized words. Whitespace is preserved as the underscore character.

\subsection{Unigram Language Model Mode}

\begin{algorithm}[Unigram LM Tokenization]
	\begin{enumerate}
		\item Start with a large seed vocabulary (all characters + frequent substrings).
		\item Define $p(x) = \sum_{s \in S(x)} p(s)$, where $S(x)$ is the set of segmentations of $x$.
		\item Optimise vocabulary by removing tokens that reduce likelihood the least.
		\item Repeat until the desired vocabulary size is reached.
	\end{enumerate}
\end{algorithm}

The unigram model defines the probability of a tokenisation as:
\begin{equation}
	p(\mathbf{x}) = \prod_{i=1}^n p(t_i), \quad \text{where } t_i \text{ are the tokens}.
\end{equation}

Training maximises the marginal likelihood $p(\text{corpus}) = \prod_{x \in \text{corpus}} \sum_{S \in \text{seg}(x)} \prod_{t \in S} p(t)$, which is computed efficiently via the forward--backward algorithm.

\section{tiktoken}

tiktoken is OpenAI's fast BPE tokenizer, implemented in Rust with Python bindings. It provides several encoding schemes:

\begin{itemize}
	\item \texttt{r50k\_base}: GPT-3 (50K vocabulary, code and text).
	\item \texttt{cl100k\_base}: GPT-3.5 and GPT-4 (100K vocabulary, extended multilingual and code support).
	\item \texttt{o200k\_base}: GPT-4o (200K vocabulary, significantly expanded multilingual and emoji coverage).
\end{itemize}

\begin{theorem}[Vocab Size Evolution]
	The progression of tokenizer vocabulary sizes reflects the need for:
	\begin{itemize}
		\item \textbf{50K (GPT-3, 2020):} sufficient for English-dominant text.
		\item \textbf{100K (GPT-3.5/4, 2023):} better multilingual coverage and code efficiency.
		\item \textbf{200K (GPT-4o, 2024):} broad multilingual parity and reduced token/byte ratio for CJK languages.
	\end{itemize}
\end{theorem}

\section{WordPiece}

WordPiece, used in BERT and its variants, differs from BPE in its merge criterion:

\begin{algorithm}[WordPiece Merge Criterion]
	Instead of merging the most frequent pair, WordPiece merges the pair $(a, b)$ that maximises the likelihood increase:
	\begin{equation}
		\Delta \mathcal{L} = \log \frac{p(b \mid a)}{p(b)} = \log \frac{\text{count}(a,b)}{\text{count}(a)\text{count}(b)}.
	\end{equation}
	This prefers pairs whose co-occurrence is higher than expected by chance.
\end{algorithm}

\section{Vocabulary Size Trade-offs}

\begin{longtable}{p{3cm}p{4cm}p{4cm}p{4cm}}
	\toprule
	\textbf{Vocab Size} & \textbf{Advantages}          & \textbf{Disadvantages}         & \textbf{Examples}   \\
	\midrule
	32K (small)         & Small embedding matrix, fast & High OOV rate, poor multiling. & Llama 1/2           \\
	50K (medium)        & Balanced for English         & CJK tokenisation inefficient   & GPT-3, GPT-2        \\
	100K (large)        & Better multilingual          & Larger embedding layer         & GPT-3.5/4, Llama 3  \\
	128K (large)        & Good coverage                & 25\%+ of params in embeddings  & Llama 3             \\
	200K (XL)           & Minimal OOV                  & Very large embedding           & GPT-4o, o200k\_base \\
	\bottomrule
\end{longtable}

For a vocabulary size $V$ and hidden dimension $d$, the embedding matrix has $V \cdot d$ parameters. At $V = 128K$ and $d = 8192$, the embedding accounts for over 1B parameters---a significant fraction of an 8B model.

\section{Multilingual Tokenization and CJK}

Languages with large character sets (Chinese, Japanese, Korean) expose weaknesses in BPE:

\begin{itemize}
	\item English text averages $\sim 1.3$ tokens/word.
	\item Chinese text averages $\sim 2.5$--$3.5$ tokens/character due to UTF-8 byte decomposition.
	\item A Chinese character (e.g., the character for ``middle'') is split into multiple bytes before BPE merging.
	\item Llama 3's 128K vocabulary was specifically designed to improve CJK efficiency.
	\item GPT-4o's 200K vocabulary dramatically reduces CJK token/byte ratios.
\end{itemize}

\begin{definition}[Token-to-Byte Ratio]
	The token-to-byte ratio measures tokenization efficiency:
	\begin{equation}
		\rho = \frac{\text{number of tokens}}{\text{number of bytes}}.
	\end{equation}
	Lower $\rho$ means more efficient encoding. For English, $\rho \approx 0.25$ (GPT-4), for Chinese, $\rho \approx 0.6$--$1.0$ depending on vocabulary size.
\end{definition}

\section{Digit Tokenization and Arithmetic Failures}

BPE tokenization of numbers causes systematic arithmetic failures:

\begin{itemize}
	\item The number ``15213'' may be tokenized as [``15'', ``213''] or [``1'', ``5213''] depending on context.
	\item Arithmetic operations that should be consistent across digit representations become non-compositional.
	\item The model must learn arithmetic independently for each digit substring, rather than leveraging place-value structure.
	\item Switching to per-digit tokenization (``1'', ``5'', ``2'', ``1'', ``3'') resolves many arithmetic failures but increases sequence length.
\end{itemize}

This is a known limitation: BPE merges frequent digit pairs (like ``00'' or ``99''), destroying the place-value structure that makes arithmetic systematic.

\section{Special Tokens}

Tokenizers reserve special tokens for structural control:

\begin{itemize}
	\item \texttt{<|begin\_of\_text|>} / \texttt{<|end\_of\_text|>}: sequence boundaries.
	\item \texttt{<|pad|>}: padding for batch processing.
	\item \texttt{<|unk|>}: unknown token (rare in byte-level tokenizers).
	\item \texttt{<|im\_start|>}, \texttt{<|im\_end|>}: chat message delimiters.
	\item \texttt{<|tool\_call|>}: structured tool invocation markers.
	\item \texttt{<|reserved|>}: reserved for task-specific tokens (e.g., special tokens for function calling in GPT-4).
\end{itemize}

The design of special tokens has become increasingly important for structured output generation and agentic workflows.

\section{Formal Semantics of BPE}

Recent work (arXiv:2309.08715) formalises BPE tokenization semantics:

\begin{theorem}[BPE Determinism]
	Byte-level BPE with a fixed merge table defines a deterministic tokenization of any string, provided no merge creates a token absent from the vocabulary.
\end{theorem}

\begin{theorem}[BPE Non-Monotonicity]
	BPE tokenization is not monotonic: adding a merge can change the tokenization of strings that did not contain the merged pair, because merges interact through shared prefixes and suffixes.
\end{theorem}

This non-monotonicity means that tokenization decisions are globally interdependent---merges cannot be applied independently.

\section{The Tokenizer Wars (2023--2024)}

The importance of tokenizer design became a debated topic:

\begin{itemize}
	\item Some argued that tokenizer choice is as important as model architecture (and that GPT-4o's quality improvements partly stem from its better tokenizer).
	\item Others countered that scaling the vocabulary is a straightforward engineering optimization, not a research breakthrough.
	\item The release of GPT-4o's \texttt{o200k\_base} encoding showed clear improvements in multilingual token/byte ratios, suggesting that prior tokenizers were suboptimal for non-English languages.
\end{itemize}

The consensus emerging by 2025: tokenizer quality matters, especially for multilingual performance and code generation, but it is a necessary baseline consideration rather than a differentiating research contribution.

\section{Exercises}

\begin{enumerate}
	\item Implement BPE merge training on a small corpus of 100 sentences.
	\item Compare the tokenization efficiency of GPT-2's 50K tokenizer vs. GPT-4o's 200K tokenizer on a Chinese Wikipedia article.
	\item Show that BPE digit tokenization makes arithmetic non-compositional.
	\item Compute the embedding parameter count for $V=128K$, $d=8192$.
	\item Write pseudocode for the SentencePiece unigram LM training loop.
	\item Explain why WordPiece's merge criterion differs from BPE's.
\end{enumerate}

% ============================================================
\chapter{Mixture of Experts}
\label{ch:moe}

\epigraph{``Sparsity is the ultimate regularizer.''}{Anonymous}

Mixture-of-Experts (MoE) layers replace dense FFN layers with multiple ``expert'' sub-networks, routing each token to a subset of experts. This decouples model capacity (total parameters) from computational cost (FLOPs per token).

\section{Top-$k$ Routing}

\begin{definition}[MoE Layer]
	An MoE layer consists of $E$ expert networks $f_1, \dots, f_E$ and a router $g$:
	\begin{equation}
		y = \sum_{i=1}^E g(x)_i \cdot f_i(x),
	\end{equation}
	where $g(x) \in \mathbb{R}^E$ is a sparse gating vector.
\end{definition}

\begin{definition}[Top-$k$ Routing]
	The router computes a score for each expert and activates only the top $k$:
	\begin{equation}
		g(x) = \text{softmax}\left(\text{TopK}\left(W_g x + \epsilon, k\right)\right),
	\end{equation}
	where $\text{TopK}(v, k)_i = v_i$ if $v_i$ is in the $k$ largest values, and $-\infty$ otherwise.
	Typically $k = 2$ for a balance of capacity and computational efficiency.
\end{definition}

The computational cost of an MoE layer with $E$ experts, top-$k$ routing, and hidden dimension $d$ is $k/E$ of the cost of a dense layer with $E \cdot d$ hidden units, while retaining $E \times$ the parameters.

\section{Token Choice vs. Expert Choice Routing}

\begin{itemize}
	\item \textbf{Token Choice Routing (TCR):} each token selects its top-$k$ experts. Simple and widely used (Switch Transformer, Mixtral). May cause load imbalance.
	\item \textbf{Expert Choice Routing (ECR):} each expert selects its top-$k$ tokens. Guarantees perfect load balance. Used in M6-T and later V-MoE (Vision MoE).
\end{itemize}

ECR addresses a fundamental flaw of TCR: under TCR, tokens may all select the same expert (catastrophic load imbalance), while other experts receive no tokens.

\begin{theorem}[ECR Load Balance Guarantee]
	With Expert Choice Routing and $T = k \cdot E$ tokens, each expert processes exactly $k$ tokens per layer, providing deterministic load balance.
\end{theorem}

\section{Load Balancing Loss}

To prevent the router from collapsing to a few experts, an auxiliary load balancing loss is added:

\begin{equation}
	\mathcal{L}_{\text{bal}} = \alpha \cdot E \cdot \sum_{i=1}^E \left( \frac{1}{T} \sum_{t=1}^T \mathbbm{1}\{\text{token } t \text{ routed to expert } i\} \right) \cdot \left( \frac{1}{T} \sum_{t=1}^T g(x_t)_i \right).
\end{equation}

This minimises the product of the fraction of tokens assigned to each expert and the average routing probability. When both are uniform ($1/E$), the loss is minimal.

\begin{remark}
	The load balancing loss trades off between routing quality (which may favour imbalanced routing) and computational efficiency (which requires balanced routing). The coefficient $\alpha$ controls this trade-off, typically set to $0.01$.
\end{remark}

\section{Fine-Grained Experts: DeepSeekMoE}

DeepSeekMoE (DeepSeek, 2024) introduces fine-grained expert splitting:

\begin{itemize}
	\item \textbf{Shared experts:} $E_s$ experts are always activated for every token.
	\item \textbf{Routed experts:} each token activates $k$ of $E_r$ routed experts via top-$k$.
\end{itemize}

The output is:
\begin{equation}
	y = \sum_{i=1}^{E_s} f_i^{\text{(shared)}}(x) + \sum_{i=1}^{E_r} g(x)_i \cdot f_i^{\text{(routed)}}(x).
\end{equation}

\begin{theorem}[Fine-Grained Expert Advantage]
	Splitting a single expert into multiple smaller experts and routing to more of them (e.g., $E=8, k=2$ vs. $E=64, k=8$) improves the knowledge-capacity trade-off: each token accesses more specialised knowledge for the same computational cost.
\end{theorem}

DeepSeek-V2 uses $E_s = 2$ shared experts and $E_r = 160$ routed experts with $k = 6$, achieving $1/\text{th}$ the inference cost of a comparable dense model.

\section{Expert Parallelism}

\begin{definition}[Expert Parallelism]
	Each expert is placed on a specific GPU (or GPU group). Tokens are routed to the GPU hosting their selected expert via all-to-all communication:
	\begin{equation}
		T_{\text{comm}} = T_{\text{all-to-all}}(B \cdot d \cdot E),
	\end{equation}
	where $B$ is the batch size and $d$ the hidden dimension.
\end{definition}

Expert parallelism introduces communication overhead: tokens must be sent to the GPU holding their expert, processed, and returned. This can become a bottleneck for small batch sizes.

\section{Grouped GEMM for MoE}

Standard MoE implementations process experts independently, launching $E$ separate matrix multiplications. Grouped GEMM fuses these into one kernel:

\begin{itemize}
	\item Tokens for all experts are packed into a contiguous input tensor.
	\item A single grouped GEMM kernel dispatches each token to its expert's weight matrix.
	\item Eliminates $E-1$ kernel launch overheads.
	\item Requires expert weight matrices arranged as $[\text{expert}, \text{hidden}, \text{intermediate}]$.
\end{itemize}

\begin{theorem}[Grouped GEMM Speedup]
	For $E$ experts and batch $B$, grouped GEMM achieves speedup:
	\begin{equation}
		S = \frac{E \cdot (T_{\text{kernel}} + T_{\text{launch}})}{T_{\text{kernel}} + E \cdot T_{\text{launch}}} \approx E \quad \text{for small batches}.
	\end{equation}
\end{theorem}

\section{PathMoE}

PathMoE (arXiv:2604.17837) introduces \emph{block-shared routers} to improve routing consistency:

\begin{itemize}
	\item Each MoE layer has a shared router $g_{\text{shared}}$ and per-expert routers $g_i$.
	\item The routing decision decomposes into a \emph{control signal} (which expert type) and a \emph{content channel} (which specific knowledge).
	\item Output: $y = g_{\text{shared}}(x) \cdot \sum_i g_i(x) \cdot f_i(x)$.
\end{itemize}

\begin{theorem}[Routing Decomposition]
	PathMoE achieves 79\% routing consistency across similar inputs, compared to $\sim 30\%$ for standard top-$k$ routing. The control signal captures the broad skill type, while the content channel specialises within that skill.
\end{theorem}

\section{Polysemantic Experts vs. Monosemantic Paths}

An important finding about MoE interpretability:

\begin{itemize}
	\item \textbf{Polysemantic experts:} individual experts activate for multiple unrelated concepts (e.g., the same expert handles both code syntax and mathematical notation).
	\item \textbf{Monosemantic paths:} the \emph{combination} of which experts are activated forms a sparse code that is more interpretable.
\end{itemize}

\begin{theorem}[Expert Superposition]
	In large MoE models, individual experts are polysemantic (they activate for multiple unrelated inputs), but the routing vector $g(x) \in \mathbb{R}^E$ acts as a sparse code that disentangles features. The effective dimensionality of the representation grows with $E$ even though each token activates only $k$ experts.
\end{theorem}

\section{Major MoE Models}

\begin{longtable}{p{3cm}p{2.5cm}p{2.5cm}p{2.5cm}p{3cm}}
	\toprule
	\textbf{Model}      & \textbf{Experts} & \textbf{Top-$k$} & \textbf{Active params} & \textbf{Notable feature}                        \\
	\midrule
	Switch Transformer  & 2048             & 1                & $\sim 7$B              & First large-scale MoE LLM, instability at $k=1$ \\
	GShard              & 2048             & 2                & $\sim 7$B              & Load balancing via capacity factor              \\
	Mixtral 8$\times$7B & 8                & 2                & 12.9B / 47B total      & First popular open MoE                          \\
	DeepSeek-V2         & 160 + 2 shared   & 6                & 21B / 236B total       & Fine-grained + shared experts                   \\
	Qwen1.5-MoE         & 64               & 8                & 2.7B / 14.3B total     & $k/E$ ratio $= 1/8$                             \\
	DBRX                & 16               & 4                & 36B / 132B total       & Fine-tuned for code                             \\
	\bottomrule
\end{longtable}

\section{Exercises}

\begin{enumerate}
	\item Derive the FLOPs comparison between a dense FFN and an MoE FFN with $E$ experts and top-$k$ routing.
	\item Show that Expert Choice Routing guarantees perfect load balance.
	\item Compute the communication volume for expert parallelism with batch $B$, hidden $d$, and $E$ experts.
	\item Explain why the load balancing loss uses the product of assignment fraction and routing probability.
	\item Compare the active-to-total parameter ratio across Mixtral, DeepSeek-V2, and DBRX.
	\item Argue why routing consistency matters for model interpretability.
\end{enumerate}

% ============================================================
\chapter{Training and Scaling Laws}
\label{ch:training}

\section{The Three-Stage Pipeline}

Modern LLM training follows a three-stage process:

\begin{enumerate}
	\item \textbf{Pre-training:} Next-token prediction on trillions of tokens. Objective:
	      \begin{equation}
		      \mathcal{L}_{\text{pretrain}} = -\frac{1}{|\mathcal{D}|} \sum_{x \in \mathcal{D}} \sum_{t} \log p_\theta(x_t \mid x_{<t}).
	      \end{equation}
	\item \textbf{Supervised fine-tuning (SFT):} Maximise likelihood of instruction--response pairs:
	      \begin{equation}
		      \mathcal{L}_{\text{SFT}} = -\frac{1}{|\mathcal{D}_{\text{SFT}}|} \sum_{(x,y) \in \mathcal{D}_{\text{SFT}}} \log p_\theta(y \mid x).
	      \end{equation}
	\item \textbf{Reinforcement learning from human feedback (RLHF):} Optimise a reward model using the PPO objective with KL penalty:
	      \begin{equation}
		      \mathcal{L}_{\text{RLHF}} = \mathbb{E}_{(x,y) \sim \pi_\theta} [r(x,y)] - \beta D_{\mathrm{KL}}(\pi_\theta \parallel \pi_{\text{SFT}}).
	      \end{equation}
\end{enumerate}

\section{Scaling Laws in Practice}

The Kaplan et al. (2020) scaling laws establish power-law relationships:
\begin{align}
	L(N) & = \left(\frac{N_c}{N}\right)^{\alpha_N} + L_\infty, \\
	L(D) & = \left(\frac{D_c}{D}\right)^{\alpha_D} + L_\infty,
\end{align}
where $\alpha_N \approx 0.076$, $\alpha_D \approx 0.103$ for typical transformer architectures.

\begin{definition}[Chinchilla Optimality]
	A model is Chinchilla-optimal when, for a given compute budget $C$, the loss $L(N, D)$ is minimised. The optimal allocation satisfies:
	\begin{equation}
		\frac{\partial L}{\partial N} \cdot \frac{dN}{dC} = \frac{\partial L}{\partial D} \cdot \frac{dD}{dC}.
	\end{equation}
	For the standard scaling law, this yields $N_\text{opt} \propto C^{0.5}$, $D_\text{opt} \propto C^{0.5}$.
\end{definition}

\begin{example}[Compute-Optimal Allocation]
	For a 70B model, Chinchilla-optimal training uses $\sim 1.4$T tokens (20 tokens per parameter). Llama 3 70B was trained on 15T tokens, making it ``overtrained'' relative to Chinchilla---a deliberate choice to maximise downstream performance.
\end{example}

\section{Learning Rate Schedules}

\begin{definition}[Cosine Decay]
	\begin{equation}
		\eta(t) = \eta_{\min} + \frac{1}{2}(\eta_{\max} - \eta_{\min})\left(1 + \cos\left(\frac{t}{T}\pi\right)\right),
	\end{equation}
	where $T$ is the total number of training steps, $\eta_{\max}$ is the peak learning rate, and $\eta_{\min}$ is the final learning rate.
\end{definition}

Modern practice (Llama 3, GPT-4):
\begin{itemize}
	\item \textbf{Warmup:} linear increase from $0$ to $\eta_{\max}$ over $2$--$5\%$ of steps.
	\item \textbf{Cosine decay:} from $\eta_{\max}$ to $\eta_{\max} \cdot 0.1$ over the remaining steps.
	\item \textbf{Peak LR:} $3\times 10^{-4}$ for 7B, $1.5\times 10^{-4}$ for 70B (inverse sqrt scaling with model size).
\end{itemize}

\section{Distributed Training Strategies}

\begin{itemize}
	\item \textbf{Data parallelism:} each GPU holds a full model replica, processes different micro-batches; gradients are all-reduced.
	\item \textbf{Tensor parallelism:} each GPU holds a slice of each layer; forward/backward passes require all-reduce on attention and FFN outputs.
	\item \textbf{Pipeline parallelism:} layers are partitioned across GPUs; micro-batches are pipelined through stages. Bubble overhead is $(P-1)/(P \cdot M)$.
	\item \textbf{ZeRO optimisation:} partitions optimizer states, gradients, and parameters across GPUs, reducing memory per GPU.
\end{itemize}

\subsection{Precision and Numerics}

\begin{itemize}
	\item \textbf{FP32:} master weights, loss scaling, normalisation.
	\item \textbf{BF16:} forward/backward pass, gradients. Larger exponent range than FP16 (no overflow for large gradients).
	\item \textbf{FP8 training (H100):} 8-bit floating point for matrix multiplications. Requires higher-precision accumulation and careful scaling.
	\item \textbf{Mixed precision:} FP32 master weights updated with FP16/BF16 gradients; loss scaling prevents underflow.
\end{itemize}

\section{Exercises}

\begin{enumerate}
	\item Compute the total FLOPs for training a 7B-parameter model on 2T tokens.
	\item Show that LoRA fine-tuning updates $\Delta W = BA$ where $B \in \mathbb{R}^{d \times r}$, $A \in \mathbb{R}^{r \times k}$ with $r \ll d,k$.
	\item Derive the memory savings of ZeRO Stage 3 over standard data parallelism.
	\item Compute the Chinchilla-optimal data budget for a 13B model.
	\item Explain why cosine decay outperforms step decay for LLM pre-training.
\end{enumerate}

% ============================================================
\chapter{Post-Training and Alignment}
\label{ch:alignment}

\epigraph{``Alignment is not a phase, it's a discipline.''}{Anonymous}

Post-training converts a pre-trained base model into a helpful, harmless, and honest assistant. This chapter covers the major alignment techniques.

\section{Supervised Fine-Tuning (SFT)}

The simplest post-training method: fine-tune on high-quality instruction-response pairs.

\begin{equation}
	\mathcal{L}_{\text{SFT}} = -\mathbb{E}_{(x,y) \sim \mathcal{D}_{\text{SFT}}}\left[\sum_{t} \log \pi_\theta(y_t \mid x, y_{<t})\right].
\end{equation}

SFT teaches the model to follow instructions and adopt a desired interaction format. It does not, however, teach the model to refuse harmful requests or to express uncertainty---these require RL-based methods.

\section{The RLHF Framework}

Reinforcement Learning from Human Feedback (RLHF) consists of three phases:

\begin{enumerate}
	\item \textbf{SFT:} initialise the policy $\pi_{\text{SFT}}$ from the pre-trained base model.
	\item \textbf{Reward modelling:} train $r_\phi(x, y)$ to predict human preferences:
	      \begin{equation}
		      \mathcal{L}_{\text{RM}} = -\mathbb{E}_{(x, y_w, y_l) \sim \mathcal{D}_{\text{RM}}}\left[\log \sigma(r_\phi(x, y_w) - r_\phi(x, y_l))\right].
	      \end{equation}
	\item \textbf{Policy optimisation:} maximise reward with KL penalty:
	      \begin{equation}
		      \max_{\pi_\theta} \mathbb{E}_{x \sim \mathcal{D}, y \sim \pi_\theta(\cdot \mid x)}\left[r_\phi(x, y)\right] - \beta D_{\mathrm{KL}}(\pi_\theta \parallel \pi_{\text{SFT}}).
	      \end{equation}
\end{enumerate}

\section{PPO (Proximal Policy Optimization)}

PPO (Schulman et al., 2017) is the workhorse of RLHF:

\begin{algorithm}[PPO-RLHF]
	For each iteration:
	\begin{enumerate}
		\item Sample trajectories $y \sim \pi_\theta(\cdot \mid x)$.
		\item Compute advantages: $A_t = \sum_{t'=t}^T \gamma^{t'-t} r_{t'} - V_\psi(x, y_{<t})$.
		\item Clipped surrogate objective:
		      \begin{equation}
			      \mathcal{L}_{\text{PPO}} = \mathbb{E}\left[\min\left(\frac{\pi_\theta(y_t \mid \dots)}{\pi_{\text{old}}(y_t \mid \dots)} A_t,\ \text{clip}\left(\frac{\pi_\theta}{\pi_{\text{old}}}, 1-\epsilon, 1+\epsilon\right) A_t\right)\right].
		      \end{equation}
		\item Update value function $V_\psi$ via MSE on returns.
		\item Add KL penalty: $\mathcal{L} = \mathcal{L}_{\text{PPO}} - \beta D_{\mathrm{KL}}(\pi_\theta \parallel \pi_{\text{SFT}})$.
	\end{enumerate}
\end{algorithm}

The critic network $V_\psi$ estimates the value function (expected return from a state), reducing variance in the policy gradient estimate.

\section{DPO (Direct Preference Optimization)}

DPO (Rafailov et al., 2023) eliminates the reward model and critic, optimising the policy directly from preferences:

\begin{equation}
	\mathcal{L}_{\text{DPO}} = -\mathbb{E}_{(x, y_w, y_l) \sim \mathcal{D}}\left[\log \sigma\left(\beta \log \frac{\pi_\theta(y_w \mid x)}{\pi_{\text{SFT}}(y_w \mid x)} - \beta \log \frac{\pi_\theta(y_l \mid x)}{\pi_{\text{SFT}}(y_l \mid x)}\right)\right].
\end{equation}

\begin{theorem}[DPO Equivalence to RLHF]
	DPO is equivalent to RLHF under the Bradley-Terry preference model: the optimal policy $\pi^*$ under the KL-constrained reward maximisation objective is identical to the policy obtained by DPO, up to the choice of $\beta$.
\end{theorem}

\begin{proof}[Sketch]
	The RLHF objective $\max_{\pi} \mathbb{E}[r(x,y)] - \beta D_{\mathrm{KL}}(\pi \parallel \pi_{\text{SFT}})$ has closed-form solution:
	\begin{equation}
		\pi^*(y \mid x) = \frac{1}{Z(x)} \pi_{\text{SFT}}(y \mid x) \exp\left(\frac{1}{\beta} r(x, y)\right).
	\end{equation}
	Solving for $r$ and substituting into the Bradley-Terry preference probability yields the DPO loss.
\end{proof}

\subsection{Likelihood Displacement Failure}

Recent work (Razin et al., 2024) identifies a failure mode of DPO:

\begin{theorem}[Likelihood Displacement]
	DPO can cause the policy to assign high probability to both preferred and dispreferred responses simultaneously, because the loss only penalises the \emph{ratio} of probabilities, not their absolute values. This ``likelihood displacement'' degrades calibration.
\end{theorem}

\section{GRPO (Group Relative Policy Optimization)}

GRPO, used by DeepSeek-R1, eliminates the critic network entirely:

\begin{algorithm}[GRPO]
	For each prompt $x$:
	\begin{enumerate}
		\item Sample $G$ responses $\{y_1, \dots, y_G\} \sim \pi_\theta(\cdot \mid x)$.
		\item Compute group-relative advantage:
		      \begin{equation}
			      A_i = \frac{r(x, y_i) - \mu_r}{\sigma_r},
		      \end{equation}
		      where $\mu_r = \frac{1}{G}\sum_j r(x, y_j)$, $\sigma_r = \text{std}(\{r(x, y_j)\})$.
		\item Policy gradient update (no critic):
		      \begin{equation}
			      \nabla_\theta \mathcal{L}_{\text{GRPO}} = \mathbb{E}\left[ \frac{1}{G}\sum_{i=1}^G \sum_{t=1}^{|y_i|} \nabla_\theta \log \pi_\theta(y_{i,t} \mid \dots) A_i - \beta \nabla_\theta D_{\mathrm{KL}}(\pi_\theta \parallel \pi_{\text{SFT}}) \right].
		      \end{equation}
	\end{enumerate}
\end{algorithm}

\begin{remark}
	GRPO replaces the learned value function with a Monte Carlo baseline computed from the group. This removes the critic network (halving memory during RL training) but introduces variance from finite group samples.
\end{remark}

\subsection{KL in the Loss vs. KL in the Reward}

GRPO places the KL penalty inside the loss rather than inside the reward:
\begin{itemize}
	\item \textbf{KL in reward:} $r'(x,y) = r(x,y) - \beta D_{\mathrm{KL}}(\pi_\theta \parallel \pi_{\text{SFT}})$. Used by PPO.
	\item \textbf{KL in loss:} $\mathcal{L} = \mathcal{L}_{\text{policy}} - \beta D_{\mathrm{KL}}(\pi_\theta \parallel \pi_{\text{SFT}})$. Used by GRPO, KTO.
\end{itemize}

When KL is in the loss, the KL penalty differentiates through the policy output, providing direct gradient signal to stay close to the reference model.

\section{RLVR (Reinforcement Learning from Verifiable Rewards)}

RLVR replaces the learned reward model with a rule-based verifier:

\begin{definition}[Rule-Based Verifier]
	For problems with objectively correct answers (math, code), the verifier $v(x, y)$ returns $1$ for correct and $0$ for incorrect, based on:
	\begin{itemize}
		\item \textbf{Math:} comparing the final numerical answer.
		\item \textbf{Code:} executing the generated code against test cases.
		\item \textbf{Logic:} checking formal constraints.
	\end{itemize}
\end{definition}

RLVR applies any policy gradient method (PPO, GRPO) with $r(x,y) = v(x,y)$:

\begin{equation}
	\max_{\pi_\theta} \mathbb{E}_{(x,y) \sim \pi_\theta}[v(x,y)] - \beta D_{\mathrm{KL}}(\pi_\theta \parallel \pi_{\text{SFT}}).
\end{equation}

\begin{remark}
	RLVR avoids reward hacking and reward model bias. DeepSeek-R1 uses RLVR for math and code, achieving significant reasoning improvements without human preference labels.
\end{remark}

\section{Comparison of RL Methods}

\begin{longtable}{p{2.5cm}p{3cm}p{3cm}p{3cm}p{3cm}}
	\toprule
	\textbf{Property} & \textbf{PPO}      & \textbf{DPO}            & \textbf{GRPO}   & \textbf{RLVR}       \\
	\midrule
	Reward model      & Required          & Not needed              & Required        & Rule-based          \\
	Critic network    & Required          & Not needed              & Not needed      & Optional            \\
	Online sampling   & Yes               & No (offline)            & Yes             & Yes                 \\
	Memory (vs. SFT)  & $2\times$         & $1\times$               & $1.5\times$     & $1\text{--}2\times$ \\
	Gradient bias     & Low               & Likelihood displacement & Group-dependent & Verifier-dependent  \\
	KL placement      & Reward            & Implicit                & Loss            & Reward or loss      \\
	Best for          & General alignment & Simple preference       & Reasoning tasks & Math/code           \\
	Used by           & GPT-4, Claude     & Llama 3, Zephyr         & DeepSeek-R1     & DeepSeek-R1, Qwen   \\
	\bottomrule
\end{longtable}

\section{Other Alignment Methods}

\begin{itemize}
	\item \textbf{KTO (Kahneman-Tversky Optimisation):} uses only binary feedback (good/bad) rather than pairwise preferences. Loss derived from prospect theory.
	\item \textbf{SimPO (Simple Preference Optimization):} replaces the reference model with sequence-level likelihood, simplifying the DPO objective.
	\item \textbf{IPO (Identity Preference Optimization):} an alternative to DPO that avoids the log-ratio formulation and provides stronger guarantees against overfitting.
	\item \textbf{ORPO (Odds Ratio Preference Optimization):} combines SFT and preference optimisation into a single loss without a separate reference model.
\end{itemize}

\subsection{K1 vs. K3 KL Estimators}

The KL divergence in RLHF can be estimated with different sample sizes:
\begin{itemize}
	\item \textbf{K1 estimator (per-token):} $\hat{D}_{\mathrm{KL}}^{(1)} = \log \frac{\pi_\theta(y_t \mid \dots)}{\pi_{\text{SFT}}(y_t \mid \dots)}$. Low variance, biased per-token.
	\item \textbf{K3 estimator (per-sequence):} $\hat{D}_{\mathrm{KL}}^{(3)} = \sum_t \log \frac{\pi_\theta(y_t \mid \dots)}{\pi_{\text{SFT}}(y_t \mid \dots)}$. Unbiased per-sequence, higher variance.
\end{itemize}

When KL is placed in the reward function, K1 is preferred because it provides a per-token KL penalty without introducing sequence-level noise. When KL is in the loss, K3 provides an unbiased estimate.

\section{Exercises}

\begin{enumerate}
	\item Derive the closed-form optimal policy for the KL-constrained RLHF objective.
	\item Show that DPO's loss function is equivalent to the RLHF objective under Bradley-Terry preferences.
	\item Compare the memory requirements of PPO vs. DPO training for a 7B model.
	\item Explain why GRPO does not need a critic network.
	\item Simulate likelihood displacement in DPO with a synthetic preference dataset.
	\item Design a rule-based verifier for a multi-step math word problem.
\end{enumerate}

% ============================================================
% Part III: Beyond Transformers
% ============================================================
\part{Beyond Transformers}

% ============================================================
\chapter{State Space Models and Alternative Architectures}
\label{ch:ssm}

\epigraph{``The quadratic cost of attention is not a law of nature; it is a design choice.''}{Anonymous}

The $O(n^2)$ scaling of attention motivates alternative architectures that achieve linear or sub-quadratic complexity while maintaining expressiveness.

\section{From RNNs to Structured State Spaces}

Recurrent neural networks (RNNs) process sequences in $O(n)$ time but suffer from vanishing gradients and limited memory. State space models (SSMs) reformulate recurrence as a continuous-time dynamical system:

\begin{equation}
	h'(t) = A h(t) + B x(t), \quad y(t) = C h(t) + D x(t),
\end{equation}
where $A \in \mathbb{R}^{d \times d}$ is the state transition matrix, $B \in \mathbb{R}^{d \times 1}$ the input projection, $C \in \mathbb{R}^{1 \times d}$ the output projection, and $D$ a skip connection.

\section{S4: Structured State Space Sequence Model}

S4 (Gu et al., 2022) introduces a structured parametrisation of $A$:

\begin{definition}[HiPPO Matrix]
	The HiPPO (High-Order Polynomial Projection Operators) matrix is:
	\begin{equation}
		A_{nk} = -\sqrt{(2n+1)(2k+1)} \cdot \begin{cases}
			1   & \text{if } n > k, \\
			n+1 & \text{if } n = k, \\
			0   & \text{if } n < k.
		\end{cases}
	\end{equation}
	This matrix stores a compressed representation of the input history as coefficients of orthogonal polynomials.
\end{definition}

S4 exploits the HiPPO matrix's normal-plus-low-rank (NPLR) structure to diagonalise the recurrence, achieving $O(n)$ computation during training via convolution.

\begin{theorem}[S4 Computational Complexity]
	S4 processes a sequence of length $n$ in $O(n)$ time and $O(n)$ memory for training, and $O(1)$ per step for generation.
\end{theorem}

\section{H3: Leveraging SSMs for Language}

H3 (Hungry Hungry Hippo; Fu et al., 2023) extends S4 by adding multiplicative gates inspired by attention:

\begin{itemize}
	\item Two SSMs in series with a gating mechanism.
	\item The first SSM produces a query-like signal; the second produces a key-like signal.
	\item Element-wise multiplication of the two SSM outputs creates a content-based interaction.
\end{itemize}

H3 achieves competitive results with transformers on language modelling while maintaining linear complexity.

\section{Mamba: S6 Selective SSM}

Mamba (Gu \& Dao, 2023) introduces the S6 layer, which makes the SSM parameters input-dependent:

\begin{equation}
	B_t = \text{Linear}_B(x_t), \quad C_t = \text{Linear}_C(x_t), \quad \Delta_t = \text{softplus}(\text{Linear}_\Delta(x_t)).
\end{equation}

\begin{definition}[Selectivity]
	A model is \emph{selective} if its state-space dynamics depend on the input content, not just position. This enables content-aware recall (e.g., ``ignore filler tokens, remember the key entity''), which was impossible in LTI SSMs.
\end{definition}

Selectivity breaks the convolution equivalence: Mamba cannot be computed via convolution because $B_t$ and $C_t$ vary per token. Instead, Mamba uses a hardware-aware parallel scan:

\begin{algorithm}[Mamba Parallel Scan]
	\begin{enumerate}
		\item Compute $\bar{A}_t = \exp(\Delta_t \cdot A)$, $\bar{B}_t = \Delta_t \cdot B_t$.
		\item Scan recurrence: $h_t = \bar{A}_t h_{t-1} + \bar{B}_t x_t$.
		\item Parallelise via associative scan: the recurrence $(h_t, \bar{A}_t, \bar{B}_t x_t)$ forms a semigroup under composition.
		\item Compute $y_t = C_t h_t$.
	\end{enumerate}
\end{algorithm}

The parallel scan runs in $O(\log n)$ steps on $O(n)$ processors, achieving practical speedups over attention for long sequences.

\section{Mamba-2: SSD Reformulation}

Mamba-2 (Dao \& Gu, 2024) reformulates the SSM as a structured state-space duality (SSD):

\begin{theorem}[SSD = Attention + SSM]
	The Mamba-2 layer computes:
	\begin{equation}
		y = (L \circ S) \cdot v,
	\end{equation}
	where $L$ is a causal mask (lower-triangular), $S = Q K^\top$ is a structured semi-separable matrix, and $\circ$ denotes element-wise multiplication.
\end{theorem}

This reveals that Mamba-2's SSM is a \emph{linear attention} variant with a structured kernel. Key innovations:
\begin{itemize}
	\item \textbf{Multi-head SSM:} multiple independent SSM heads, matching attention's multi-head structure.
	\item \textbf{Head dimension:} $d_{\text{head}} = 64$ or 128, matching typical attention configurations.
	\item \textbf{State dimension:} $N = 64$--$256$ per head, larger than Mamba-1's $N=16$.
\end{itemize}

The SSD formulation enables efficient training via matrix multiplication ($O(n^2)$ on GPU, matching attention), while retaining $O(1)$ per-step generation.

\section{Mamba-3: Complex States and MIMO}

Mamba-3 (2025) introduces three major advances:

\begin{itemize}
	\item \textbf{Exponential-trapezoidal discretisation:} higher-order integration of the continuous SSM reduces approximation error.
	\item \textbf{Complex-valued states:} $A \in \mathbb{C}^{d \times d}$ with complex eigenvalues enables oscillatory dynamics and richer temporal patterns.
	\item \textbf{MIMO (multi-input multi-output):} each SSM processes multiple input and output channels simultaneously, increasing expressiveness per parameter.
\end{itemize}

\begin{theorem}[Mamba-3 Expressiveness]
	With complex-valued states, Mamba-3 can represent any linear dynamical system of dimension $2d$, whereas real-valued Mamba-2 requires $2d$ real dimensions.
\end{theorem}

\section{Hybrid Architectures}

Several models combine SSMs with attention for the best of both worlds:

\begin{itemize}
	\item \textbf{Jamba (AI21, 2024):} interleaves Mamba layers with attention layers and MoE FFNs. Attention layers handle long-range recall; Mamba layers handle efficient local processing.
	\item \textbf{TransMamba:} shares query/key/value projections between the attention and SSM paths, using a cross-block residual (CBx) connection. Params: 91\% shared, 9\% unique.
	\item \textbf{RWKV:} linear attention reformulated as a recurrent neural network with a time-dependent decay factor. $O(n)$ training via WKV computation.
	\item \textbf{RetNet (Sun et al., 2023):} retention mechanism with three computation modes---parallel (training), recurrent (generation), chunkwise (efficient decoding).
\end{itemize}

\section{MoE-SSM Hybrids}

\begin{itemize}
	\item \textbf{Swimba (2025):} replaces FFN experts in MoE with SSM experts. The router selects which SSM to apply per token. Single-pass recurrence (no KV cache) enables extreme throughput.
	\item \textbf{Nemotron-H (NVIDIA, 2025):} hybrid MoE with both attention and SSM experts, using hardware-aware scheduling to balance workload across experts.
\end{itemize}

\section{Comparison with Transformers}

\begin{longtable}{p{3cm}p{3cm}p{3cm}p{3cm}}
	\toprule
	\textbf{Property}       & \textbf{Transformer} & \textbf{Mamba-2}    & \textbf{Jamba}  \\
	\midrule
	Training complexity     & $O(n^2 d)$           & $O(n d N)$          & Mixed           \\
	Generation complexity   & $O(n d)$ per step    & $O(1)$ per step     & $O(1)$ per step \\
	Memory (long context)   & $O(n^2)$             & $O(n d)$            & $O(n d)$        \\
	Content-based recall    & Yes (attention)      & Yes (selective)     & Yes             \\
	Long-range dependencies & Unlimited            & State-limited ($N$) & Hybrid          \\
	Throughput (long seq.)  & Degrades             & Stable              & Near-stable     \\
	\bottomrule
\end{longtable}

\section{Exercises}

\begin{enumerate}
	\item Show that S4 can be computed via convolution when $A$, $B$, $C$ are time-invariant.
	\item Derive the associative scan for the Mamba recurrence $h_t = \bar{A}_t h_{t-1} + \bar{B}_t x_t$.
	\item Explain why selectivity breaks the convolution equivalence in Mamba.
	\item Prove that the SSD formulation in Mamba-2 is equivalent to linear attention with a structured kernel.
	\item Compare the generation-time memory of a transformer (KV cache) vs.\ Mamba (state vector).
	\item Design a hybrid architecture that interleaves 1 attention layer every 4 Mamba layers. Justify the ratio.
\end{enumerate}

% ============================================================
% Part IV: Inference
% ============================================================
\part{Inference}

% ============================================================
\chapter{Autoregressive Inference}
\label{ch:inference}

\section{Decoding Strategies}

\begin{definition}[Greedy Decoding]
	At each step $t$, select
	\begin{equation}
		x_t^* = \arg\max_{x \in \mathcal{V}} p_\theta(x \mid x_{<t}).
	\end{equation}
	Greedy decoding is deterministic but may miss high-probability continuations due to locally optimal choices.
\end{definition}

\begin{definition}[Beam Search]
	Maintain a set $\mathcal{B}_t$ of $k$ partial hypotheses. At each step, expand each hypothesis with the top-$k$ tokens and retain the $k$ highest-scoring sequences:
	\begin{equation}
		\mathcal{B}_{t+1} = \text{top-}k \left( \bigcup_{y \in \mathcal{B}_t} \bigcup_{x \in \mathcal{V}} \{(y, x)\} \text{ with score } \log p_\theta(y, x) \right).
	\end{equation}
\end{definition}

\begin{definition}[Top-$k$ Sampling]
	Sample from the $k$ highest-probability tokens:
	\begin{equation}
		p'(x \mid x_{<t}) = \begin{cases}
			p(x \mid x_{<t}) / \sum_{i=1}^k p_{(i)} & \text{if } x \in \mathcal{V}^{(k)}, \\
			0                                       & \text{otherwise},
		\end{cases}
	\end{equation}
	where $\mathcal{V}^{(k)}$ is the set of $k$ tokens with highest probability.
\end{definition}

\begin{definition}[Top-$p$ (Nucleus) Sampling]
	Sample from the smallest set $\mathcal{V}^{(p)}$ such that
	\begin{equation}
		\sum_{x \in \mathcal{V}^{(p)}} p(x \mid x_{<t}) \geq p.
	\end{equation}
	Nucleus sampling adaptively truncates the tail of the distribution.
\end{definition}

\section{Temperature Sampling}

Temperature $\tau$ rescales the logits before the softmax:
\begin{equation}
	p_\tau(x \mid x_{<t}) = \frac{\exp(\ell_t / \tau)}{\sum_{x'} \exp(\ell_{x'} / \tau)}.
\end{equation}

As $\tau \to 0$, sampling approaches greedy decoding. As $\tau \to \infty$, sampling approaches uniform.

\begin{theorem}[Temperature and Entropy]
	The entropy of the temperature-scaled distribution satisfies:
	\begin{equation}
		H(p_\tau) = \log|\mathcal{V}| - \frac{1}{\tau} \sum_x p(x) \log p(x) + O(\tau^{-2}) \quad \text{as } \tau \to \infty.
	\end{equation}
	Lower temperatures concentrate probability mass, reducing diversity.
\end{theorem}

\section{The KV Cache}

During autoregressive generation, the key and value tensors for all previous tokens are cached to avoid recomputation:
\begin{align}
	K_{\leq t} & = [K_{\leq t-1}; K_t] \in \mathbb{R}^{t \times d_k}, \\
	V_{\leq t} & = [V_{\leq t-1}; V_t] \in \mathbb{R}^{t \times d_v}.
\end{align}

The KV cache has memory cost:
\begin{equation}
	M_{\text{KV}} = 2 \cdot L \cdot n_{kv} \cdot t \cdot d_h \cdot p \quad \text{bytes}.
\end{equation}

For a 70B model with $t = 4096$, $L=80$, $n_{kv}=8$, $d_h=128$, $p=2$ (FP16): $M \approx 1.34$ GB per sequence.

\section{Parallel Decoding}

Several techniques generate multiple tokens per step:

\begin{itemize}
	\item \textbf{Medusa (Cai et al., 2024):} add $k$ parallel heads that predict the next $k$ tokens simultaneously. Each head predicts one future position. Use tree attention to process all hypotheses.
	\item \textbf{Lookahead Decoding (Fu et al., 2024):} use Jacobi iteration to generate $k$ tokens per step. Each step, the model predicts all $k$ positions based on current guesses, then updates.
	\item \textbf{Eagle (Li et al., 2024):} a lightweight draft model predicts future tokens using both the current hidden state and previously predicted tokens.
\end{itemize}

These methods achieve 2--4$\times$ speedup in wall-clock time at the cost of increased FLOPs.

\section{Exercises}

\begin{enumerate}
	\item Compare greedy decoding and beam search on a sequence where the locally optimal choice leads to a dead end.
	\item Compute the KV cache memory for Llama 3 70B at $t = 8192$.
	\item Show that temperature $\tau = 1$ recovers the original distribution.
	\item Design a Medusa head architecture for 4-token lookahead.
	\item Explain why nucleus sampling is preferred over top-$k$ for diverse generation.
\end{enumerate}

% ============================================================
\chapter{Inference Optimisation}
\label{ch:inference-opt}

\section{Quantisation}

\begin{definition}[Uniform Quantisation]
	Map a weight $w \in [\alpha, \beta]$ to $b$-bits:
	\begin{equation}
		\hat{w} = \Delta \cdot \text{clamp}\left(\left\lfloor \frac{w - \alpha}{\Delta} \right\rceil, 0, 2^b - 1\right) + \alpha,
	\end{equation}
	where $\Delta = (\beta - \alpha) / (2^b - 1)$.
\end{definition}

\subsection{GPTQ}

Optimal Brain Quantiser (Frantar et al., 2023): iteratively quantise weights while compensating the remainder using the inverse Hessian:
\begin{equation}
	\delta w_q = -\frac{w_q}{\text{diag}(H_{qq}^{-1})} \cdot H_{:,q}^{-1},
\end{equation}
where $H = 2XX^\top$ is the Hessian of the MSE loss.

\subsection{AWQ}

Activation-aware weight quantisation (Lin et al., 2023): scale salient weights (those corresponding to large activations) before quantisation to preserve outlier channels:
\begin{equation}
	\tilde{w} = w \cdot s, \quad \text{Quantise}(\tilde{w}), \quad \hat{w} = \text{Dequantise}(\hat{\tilde{w}}) \cdot s^{-1}.
\end{equation}
The per-channel scale $s$ is chosen to minimise MSE under a small calibration set.

\subsection{GGUF and NF4}

For CPU inference, llama.cpp uses the GGUF format with various quantisation levels:
\begin{itemize}
	\item \textbf{GGUF Q4\_K\_M:} 4-bit quantisation with importance-based mixed precision (some weights get 5 or 6 bits).
	\item \textbf{NF4 (QLoRA):} normalised float 4-bit, where quantisation levels match a standard normal distribution: $\text{NF4} = \{ -1.0, -0.85, -0.7, \dots, 0.85, 1.0 \}$.
\end{itemize}

\subsection{FP8 Inference (H100)}

The H100 GPU supports FP8 matrix multiplication natively:
\begin{itemize}
	\item E4M3 (4 exponent bits, 3 mantissa): range $\pm 448$, fine-grained for weights.
	\item E5M2 (5 exponent bits, 2 mantissa): wider range $\pm 57344$, coarser granularity, suitable for activations.
	\item FP8 inference achieves 2$\times$ throughput vs. FP16 with minimal quality loss ($< 0.5\%$ perplexity increase).
\end{itemize}

\section{Speculative Decoding}

Speculative decoding (Leviathan et al., 2022; Chen et al., 2023) accelerates generation by using a small draft model $M_q$ to propose $K$ tokens, which are then verified by the target model $M_p$ in parallel.

\begin{algorithm}[Speculative Decoding]
	\begin{enumerate}
		\item Draft model $M_q$ generates $K$ tokens autoregressively: $x_{t+1}, \dots, x_{t+K} \sim M_q(\cdot \mid x_{\leq t})$.
		\item Target model $M_p$ computes logits for all $K+1$ prefixes in parallel.
		\item For each position $i = 1, \dots, K$, accept $x_{t+i}$ with probability $\min(1, p(x_{t+i}) / q(x_{t+i}))$.
		\item Reject at the first rejection; resample from the residual distribution $\propto \max(0, p - q)$.
	\end{enumerate}
\end{algorithm}

The expected speedup is:
\begin{equation}
	\mathbb{E}[\text{speedup}] = \frac{\sum_{k=1}^K k \cdot P(\text{accept } k)}{\text{latency overhead}}.
\end{equation}

\begin{theorem}[Acceptance Rate]
	The expected number of accepted tokens is:
	\begin{equation}
		\mathbb{E}[K_{\text{acc}}] = \sum_{k=1}^K P(\text{accept } k) = \sum_{k=1}^K \prod_{i=1}^k \min\left(1, \frac{p(x_{t+i})}{q(x_{t+i})}\right).
	\end{equation}
	When $M_q \approx M_p$, $\mathbb{E}[K_{\text{acc}}] \approx K$.
\end{theorem}

\section{Flash Attention}

Flash Attention (Dao et al., 2022) computes exact attention without materialising the $N \times N$ matrix, using tiling and recomputation:

\begin{enumerate}
	\item Tile $Q$, $K$, $V$ into blocks that fit in SRAM.
	\item For each block, compute the partial softmax and update the output incrementally using online softmax.
	\item Recompute the attention matrix during the backward pass (instead of storing it) to trade FLOPs for memory bandwidth.
\end{enumerate}

This reduces memory from $O(N^2)$ to $O(N)$ and achieves 2--4$\times$ speedup on GPUs.

\begin{theorem}[Flash Attention IO Complexity]
	Flash Attention achieves $O(N^2 d^2 / M)$ HBM accesses, where $M$ is the SRAM size, compared to $O(N^2 d + N d^2)$ for standard attention. For $M = 108$KB (A100), this is a $5$--$10\times$ reduction in HBM traffic.
\end{theorem}

\section{Continuous Batching and PagedAttention}

\begin{itemize}
	\item \textbf{Continuous batching:} instead of waiting for all sequences in a batch to complete, new sequences are admitted as soon as existing ones finish. This maximises GPU utilisation.
	\item \textbf{PagedAttention (vLLM):} KV cache blocks are stored in non-contiguous physical memory pages, managed like virtual memory. This eliminates internal fragmentation and enables copy-on-write for parallel sampling.
\end{itemize}

\section{Exercises}

\begin{enumerate}
	\item Show that 4-bit quantisation compresses model weights by a factor of $8\times$ relative to FP16.
	\item Derive the expected speedup of speculative decoding when the acceptance rate is $\alpha$.
	\item Explain why Flash Attention reduces memory from $O(N^2)$ to $O(N)$.
	\item Compare continuous batching with static batching in terms of throughput and latency.
	\item Compute the FP8 memory savings for a 70B model vs. FP16.
	\item Show that NF4 quantisation preserves the expected value of standard-normal-distributed weights.
\end{enumerate}

% ============================================================
\chapter{Long-Context Techniques}
\label{ch:long-context}

\epigraph{``Context length is not a ceiling, it is a floor.''}{Anonymous}

Standard RoPE has limited ability to generalise beyond the training sequence length. This chapter covers techniques to extend the context window of pre-trained LLMs.

\section{Position Interpolation (PI)}

PI (Chen et al., 2023) rescales position indices to fit within the trained range:

\begin{equation}
	\theta'_j = \theta_j \cdot \frac{L_{\text{train}}}{L_{\text{target}}},
\end{equation}
where $L_{\text{train}}$ is the training context length and $L_{\text{target}}$ is the desired length.

\begin{definition}[Position Interpolation]
	For a target length $L > L_{\text{train}}$, the position $m$ is mapped to $m' = m \cdot (L_{\text{train}} / L)$. All RoPE frequencies remain the same but positions are compressed.
\end{definition}

PI requires fine-tuning on the extended length but causes resolution loss for nearby tokens.

\section{NTK-Aware Scaling}

NTK-aware scaling (bloc97, 2023) scales high-frequency components less and low-frequency components more, motivated by the Neural Tangent Kernel:

\begin{equation}
	\theta'_j = \theta_j \cdot \alpha^{2j/d},
\end{equation}
where $\alpha = (L_{\text{target}} / L_{\text{train}})^{d/(d-2)}$.

\begin{remark}
	NTK-aware scaling preserves high-frequency information (local token interactions) while extending the range of low-frequency components (long-range dependencies). It requires minimal or no fine-tuning.
\end{remark}

\section{YaRN (Yet another RoPE extensioN)}

YaRN (Peng et al., 2023) combines PI and NTK-aware scaling with a ramp function:

\begin{align}
	\theta'_j       & =
	\begin{cases}
		\theta_j \cdot s    & \text{if } j < r \text{ (low freq.)},      \\
		\theta_j            & \text{if } j > r + w \text{ (high freq.)}, \\
		\text{interpolated} & \text{otherwise},
	\end{cases} \\
	\text{where } s & = \frac{L_{\text{target}}}{L_{\text{train}}}.
\end{align}

Key parameters:
\begin{itemize}
	\item \textbf{Ramp length $w$:} width of the transition region (typically $w = 8$).
	\item \textbf{Ramp start $r$:} dimension below which full interpolation is applied.
	\item \textbf{Temperature factor $t$:} rescales attention logits: $t = \sqrt{\log(L_{\text{target}} / L_{\text{train}})}$.
\end{itemize}

\begin{theorem}[YaRN Frequency Groups]
	YaRN partitions dimensions into three groups:
	\begin{enumerate}
		\item \textbf{High-frequency:} $j > r + w$, keep original $\theta_j$ (no interpolation).
		\item \textbf{Low-frequency:} $j < r$, apply full position interpolation.
		\item \textbf{Ramp region:} $r \leq j \leq r + w$, smoothly transition between the two.
	\end{enumerate}
	This preserves local information while extending range.
\end{theorem}

\section{Dynamic NTK}

Dynamic NTK applies NTK-aware scaling at inference time without fine-tuning:
\begin{equation}
	\alpha = \max\left(1, \frac{L_{\text{current}}}{L_{\text{train}}}\right),
\end{equation}
where $L_{\text{current}}$ is the actual sequence length at each inference step. The scaling factor changes dynamically as generation progresses.

\section{LongRoPE}

LongRoPE (2024) uses evolutionary search to find per-dimension position interpolation factors:

\begin{equation}
	\theta_j' = \frac{\theta_j}{s_j},
\end{equation}
where $s_j$ is the per-dimension scale, found via:
\begin{equation}
	\min_{s_1, \dots, s_{d/2}} \text{PPL}(\text{model with scales } s_j),
	\quad \text{s.t. } s_j \text{ monotonically non-decreasing in } j.
\end{equation}

\begin{definition}[Monotonic Constraint]
	The scale factors $s_j$ must be monotonically non-decreasing: $s_1 \leq s_2 \leq \dots \leq s_{d/2}$. This ensures that lower-frequency dimensions (larger $j$) are interpolated more, consistent with the NTK intuition.
\end{definition}

LongRoPE extends context to 256K with minimal perplexity degradation and no fine-tuning beyond a small calibration set.

\section{LongRoPE2}

LongRoPE2 (2025) refines the evolutionary search with a needle-driven objective:

\begin{equation}
	\min_{s} \text{PPL}_{\text{needle}}(s) = \text{PPL}(\text{needle-in-haystack with scales } s).
\end{equation}

Instead of perplexity on long documents, LongRoPE2 optimises for the needle-in-haystack retrieval success rate, directly targeting the use case that long-context models must support.

\section{AlphaRoPE}

AlphaRoPE uses power-function scaling as a compromise between PI and NTK:

\begin{equation}
	\theta_j' = \theta_j \cdot \left(\frac{L_{\text{target}}}{L_{\text{train}}}\right)^{\alpha_j},
	\quad \alpha_j = 1 - \frac{2j}{d}.
\end{equation}

High-frequency dimensions ($j \approx 0$) get $\alpha_j \approx 1$ (full PI scaling), while low-frequency dimensions ($j \approx d/2$) get $\alpha_j \approx 0$ (no scaling). This smooth interpolation avoids the hard ramp boundary of YaRN.

\section{CoPE (Contextual Position Encoding)}

CoPE (2025) applies soft clipping to low-frequency components:

\begin{equation}
	\theta_j' = \theta_j \cdot \text{clip}\left(\frac{L_{\text{target}}}{L_{\text{train}}}, 1, \tau_j\right),
\end{equation}
where $\tau_j$ is a per-dimension clipping threshold learned or defined as $\tau_j = 1 + (c-1) \cdot \frac{2j}{d}$.

CoPE achieves $2\times$ perplexity improvement over PI on long-context benchmarks without fine-tuning.

\section{RoPE Out-of-Distribution Theory}

Recent theoretical work unifies these techniques:

\begin{theorem}[RoPE OOD Decomposition]
	The OOD error from length extrapolation decomposes into:
	\begin{equation}
		\mathcal{E}_{\text{ext}} = \mathcal{E}_{\text{freq}} + \mathcal{E}_{\text{pos}} + \mathcal{E}_{\text{atten}},
	\end{equation}
	where:
	\begin{itemize}
		\item $\mathcal{E}_{\text{freq}}$: frequency mismatch between trained and novel positions.
		\item $\mathcal{E}_{\text{pos}}$: absolute position encoding drift beyond trained range.
		\item $\mathcal{E}_{\text{atten}}$: attention entropy collapse at unseen positions.
	\end{itemize}
	NTK-aware scaling reduces $\mathcal{E}_{\text{freq}}$, YaRN's temperature factor reduces $\mathcal{E}_{\text{atten}}$, and LongRoPE's per-dimension scales reduce $\mathcal{E}_{\text{pos}}$.
\end{theorem}

\section{Comparison}

\begin{longtable}{p{3cm}p{3cm}p{3cm}p{3cm}p{2.5cm}}
	\toprule
	\textbf{Method} & \textbf{Max length}         & \textbf{Fine-tune needed?} & \textbf{Key idea}    & \textbf{PPL @ 128K} \\
	\midrule
	None            & $L_{\text{train}}$          & No                         & --                   & Degrades            \\
	PI              & $2\times L_{\text{train}}$  & Yes                        & Linear rescaling     & $+1.5$              \\
	NTK-aware       & $4\times L_{\text{train}}$  & No                         & Frequency-dependent  & $+0.8$              \\
	YaRN            & $16\times L_{\text{train}}$ & No                         & Ramp + temperature   & $+0.3$              \\
	Dynamic NTK     & $8\times L_{\text{train}}$  & No                         & Evolving scale       & $+0.6$              \\
	LongRoPE        & $32\times L_{\text{train}}$ & Minimal                    & Per-dimension search & $+0.1$              \\
	AlphaRoPE       & $16\times L_{\text{train}}$ & No                         & Power-function       & $+0.4$              \\
	CoPE            & $16\times L_{\text{train}}$ & No                         & Soft clipping        & $+0.2$              \\
	\bottomrule
\end{longtable}

\section{Exercises}

\begin{enumerate}
	\item Derive the position interpolation formula for RoPE.
	\item Explain why NTK-aware scaling preserves high-frequency components.
	\item Implement the YaRN ramp function in pseudocode.
	\item Show that LongRoPE's monotonic constraint is consistent with NTK intuition.
	\item Compare the fine-tuning requirements of PI vs. YaRN vs. LongRoPE.
	\item Propose a new RoPE extension method combining AlphaRoPE and CoPE.
\end{enumerate}

% ============================================================
% Part V: Engineering
% ============================================================
\part{Engineering}

% ============================================================
\chapter{GPU Architecture for LLMs}
\label{ch:gpu}

\section{GPU Memory Hierarchy}

Modern GPUs (NVIDIA Hopper, Ada Lovelace) have a hierarchical memory system:

\begin{itemize}
	\item \textbf{Global memory (HBM3/HBM2e):} 40--144\,GB, bandwidth 2--3.35\,TB/s. High latency ($\sim$600 cycles).
	\item \textbf{L2 cache:} 40--60\,MB on-chip. Shared across SMs.
	\item \textbf{Shared memory / SRAM:} 128--256\,KB per SM. Low latency ($\sim$20 cycles), explicitly managed.
	\item \textbf{Registers:} 65536 per SM. Zero-latency access.
\end{itemize}

\section{Compute vs. Memory Bound Analysis}

An operation is:
\begin{itemize}
	\item \textbf{Compute-bound} if arithmetic intensity $> \frac{\text{peak FLOPS}}{\text{memory bandwidth}}$ (e.g., FFN layers).
	\item \textbf{Memory-bound} otherwise (e.g., attention with KV cache, embedding lookup).
\end{itemize}

For an A100 (312\,TFLOPS FP16, 2\,TB/s HBM bandwidth), the threshold is $I_{\text{crit}} = 312 / 2 = 156$\,FLOPs/byte. FFN layers: $I \approx 200$\,FLOPs/byte (compute-bound). Attention: $I \approx 20$\,FLOPs/byte (memory-bound).

\section{Tensor Cores and Mixed Precision}

Tensor Cores perform $4 \times 4$ matrix multiply-accumulate in one cycle:
\begin{equation}
	D = A \cdot B + C,
\end{equation}
where $A$ and $B$ are FP16/BF16/INT8 and $D$ and $C$ are FP16/INT32.

Mixed-precision training uses FP16/BF16 for forward/backward with FP32 master weights for numerical stability.

\subsection{Hopper-Specific Features}

NVIDIA H100 introduces:
\begin{itemize}
	\item \textbf{FP8 Tensor Cores:} 2$\times$ throughput vs. FP16, with E4M3/E5M2 formats.
	\item \textbf{Transformer Engine:} automatic FP8/FP16 selection per layer based on activation statistics.
	\item \textbf{DPX instructions:} accelerate dynamic programming (Smith-Waterman, Floyd-Warshall) used in some sequence alignment tasks.
	\item \textbf{Aynchronous transactions:} overlapping data movement with computation via thread block clusters.
\end{itemize}

\section{Parallelism Strategies}

\begin{itemize}
	\item \textbf{Tensor parallelism (TP):} each layer's parameters are sharded across GPUs. Requires all-reduce per sub-layer. Optimal when intra-node NVLink bandwidth is high.
	\item \textbf{Pipeline parallelism (PP):} layers are partitioned into stages. Bubble overhead is $\frac{(P-1) \cdot (M-1)}{P \cdot M}$ where $P$ is the number of pipeline stages and $M$ is the number of micro-batches.
	\item \textbf{Sequence parallelism:} the sequence dimension is split across GPUs. Each GPU processes a segment of the sequence.
	\item \textbf{Expert parallelism (MoE):} each expert is placed on a different GPU; tokens are routed via all-to-all communication.
\end{itemize}

\subsection{Interconnects}

\begin{itemize}
	\item \textbf{NVLink (H100):} 900 GB/s bidirectional per GPU, direct GPU-to-GPU.
	\item \textbf{InfiniBand (HDR200):} 400 Gb/s per link, used for inter-node communication.
	\item \textbf{PCIe Gen5:} 128 GB/s per direction, used for host-to-GPU transfer.
\end{itemize}

\section{Kernel Fusion}

Launching many small GPU kernels incurs overhead. Kernel fusion combines operations:
\begin{itemize}
	\item \textbf{Fused attention:} combine $QK^\top$, softmax, dropout, $SV$ into one kernel.
	\item \textbf{Fused layer norm + residual + quantisation:} one kernel for all three.
	\item \textbf{Flash Attention} (see \S10.3) is the canonical fused attention kernel.
\end{itemize}

\section{Exercises}

\begin{enumerate}
	\item Compute the arithmetic intensity of a single transformer FFN layer.
	\item Explain why memory-bound operations benefit more from kernel fusion.
	\item Compare the communication costs of TP and PP for a 70B model on 8 GPUs.
	\item Calculate the HBM bandwidth utilisation for an H100 running FP8 attention at batch size 1, sequence length 8192.
	\item Show that the roofline model predicts compute-bound behaviour for FFN layers but memory-bound for attention.
\end{enumerate}

% ============================================================
\chapter{Software Engineering for LLMs}
\label{ch:se}

\section{Prompt Engineering Patterns}

\begin{definition}[Chain-of-Thought (CoT)]
	Elicit intermediate reasoning before the final answer:
	\begin{lstlisting}
  Q: {question}
  A: Let's think step by step.
  \end{lstlisting}
\end{definition}

\begin{definition}[Few-Shot Prompting]
	Provide $k$ exemplars of the task in the prompt:
	\begin{lstlisting}
  Input:  x_1  ->  Output:  y_1
  Input:  x_2  ->  Output:  y_2
  ...
  Input:  x_k  ->  Output:  y_k
  Input:  x    ->  Output:
  \end{lstlisting}
\end{definition}

\begin{definition}[Structured Output Parsing]
	Request machine-parseable output formats:
	\begin{lstlisting}
  Output in JSON:
  {
    "answer": "...",
    "reasoning": "...",
    "confidence": 0.95
  }
  \end{lstlisting}
\end{definition}

\section{Retrieval-Augmented Generation (RAG)}

RAG augments LLM input with retrieved documents:
\begin{equation}
	p(y \mid x) = \sum_{d \in \text{Retrieve}(x)} p_{\text{LLM}}(y \mid x, d) p_{\text{retriever}}(d \mid x).
\end{equation}

The canonical RAG pipeline:
\begin{enumerate}
	\item \textbf{Indexing:} chunk documents into passages; embed with a dense retriever (e.g., BERT, E5).
	\item \textbf{Retrieval:} given query $q$, find top-$k$ passages by cosine similarity: $\text{top-}k \cos(E(q), E(D))$.
	\item \textbf{Generation:} concatenate query and retrieved passages as LLM input.
\end{enumerate}

\subsection{Advanced RAG Patterns}

\begin{itemize}
	\item \textbf{Hierarchical retrieval:} first retrieve documents, then passages within documents.
	\item \textbf{Query rewriting:} use the LLM to rewrite the query before retrieval.
	\item \textbf{HyDE (Hypothetical Document Embedding):} generate a hypothetical answer, then retrieve based on that answer's embedding.
	\item \textbf{Iterative retrieval:} after each generation step, retrieve additional context.
	\item \textbf{Self-RAG:} the LLM decides when to retrieve, what to retrieve, and whether to use retrieved information.
\end{itemize}

\section{Agent Architectures}

LLM agents extend the model with tools and structured reasoning:
\begin{lstlisting}
  def agent_loop(query):
      messages = [system_prompt, user_message(query)]
      while not done:
          response = llm(messages)
          action = parse_action(response)
          if action.type == "call_tool":
              result = execute_tool(action.name, action.args)
              messages.append({"role": "tool", "content": result})
          elif action.type == "final_answer":
              return action.content
\end{lstlisting}

\section{Testing and Evaluation}

\begin{itemize}
	\item \textbf{Unit testing:} test individual prompt templates, parsing logic, and tool integrations.
	\item \textbf{Regression testing:} maintain a set of canonical inputs and expected outputs; measure pass rate across model versions.
	\item \textbf{Evaluation suites:} MMLU (knowledge), HumanEval (code), MT-Bench (conversation), HELM (holistic).
	\item \textbf{Hallucination detection:} measure token-level perplexity, semantic entropy, or self-consistency across multiple samples.
\end{itemize}

\section{Production Deployment Patterns}

\begin{itemize}
	\item \textbf{Serve with an inference engine:} vLLM, TGI, TensorRT-LLM.
	\item \textbf{Load balancing:} round-robin or least-connections across model replicas.
	\item \textbf{Caching:} KV cache reuse for common prefixes; semantic caching for identical queries.
	\item \textbf{Monitoring:} latency, throughput, TTFT (time to first token), ITL (inter-token latency).
	\item \textbf{Guardrailing:} input/output content filtering, PII detection, topic blocking.
\end{itemize}

\section{Exercises}

\begin{enumerate}
	\item Implement a RAG pipeline pseudocode with chunking, embedding, retrieval, and generation.
	\item Design an agent that can use a calculator, a web search tool, and a database query tool.
	\item Propose a testing strategy for an LLM-based code generation system.
	\item Compare the latency impact of adding a retrieval step to an LLM call.
\end{enumerate}

% ============================================================
% Part VI: Practice
% ============================================================
\part{Practice}

% ============================================================
\chapter{Libraries and Frameworks}
\label{ch:libraries}

\section{HuggingFace Transformers}

The de facto standard library for LLMs:
\begin{lstlisting}
  from transformers import AutoModelForCausalLM, AutoTokenizer

  model = AutoModelForCausalLM.from_pretrained(
      "meta-llama/Llama-3.1-8B-Instruct",
      torch_dtype=torch.bfloat16,
      device_map="auto",
      attn_implementation="flash_attention_2",
  )
  tokenizer = AutoTokenizer.from_pretrained("meta-llama/Llama-3.1-8B-Instruct")
\end{lstlisting}

Key features: model hub, pipeline API, trainer, quantisation (bitsandbytes), PEFT (LoRA).

\section{vLLM}

High-throughput serving with PagedAttention:
\begin{lstlisting}
  from vllm import LLM, SamplingParams

  llm = LLM("meta-llama/Llama-3.1-8B-Instruct")
  params = SamplingParams(temperature=0.7, top_p=0.9, max_tokens=512)
  outputs = llm.generate(prompts, params)
\end{lstlisting}

Key features: continuous batching, PagedAttention, prefix caching, tensor parallelism, AWQ/GPTQ quantisation, OpenAI-compatible API server.

\section{TensorRT-LLM}

NVIDIA's optimised inference framework:
\begin{lstlisting}
  # Build engine
  trtllm-build --model_dir llama-8b \
    --dtype bfloat16 \
    --max_batch_size 64 \
    --max_input_len 4096 \
    --max_output_len 1024

  # Run server
  python scripts/launch_triton_server.py \
    --model_repository=/models/llama-8b \
    --world_size=4
\end{lstlisting}

Key features: FP8 quantisation, in-flight batching, multi-GPU tensor parallelism, CUDA graph optimisation, speculative decoding.

\section{llama.cpp}

CPU-first inference with GGUF quantised models:
\begin{lstlisting}
  ./llama-cli -m llama-8b.Q4_K_M.gguf \
    -p "Once upon a time" \
    -n 256 -t 8
\end{lstlisting}

Key features: GGUF quantisation (2--8 bit), CPU/GPU hybrid, Metal/CUDA/Vulkan backends, no Python dependency.

\section{LangChain}

Framework for LLM application development:
\begin{lstlisting}
  from langchain.chains import RetrievalQA
  from langchain_community.vectorstores import FAISS
  from langchain_huggingface import HuggingFaceEmbeddings

  retriever = FAISS.from_texts(docs, HuggingFaceEmbeddings()).as_retriever()
  qa = RetrievalQA.from_llm(llm, retriever=retriever)
  result = qa.invoke("What is the capital of France?")
\end{lstlisting}

Key features: modular chains, agent frameworks, tool integration, memory, and streaming.

\section{LlamaIndex}

Data framework for RAG applications:
\begin{lstlisting}
from llama_index.core import VectorStoreIndex, SimpleDirectoryReader

documents = SimpleDirectoryReader("data/").load_data()
index = VectorStoreIndex.from_documents(documents)
query_engine = index.as_query_engine()
  response = query_engine.query("What does the document say about X?")
\end{lstlisting}

Key features: advanced indexing strategies (tree, keyword, hybrid), query routing, document agents, structured data extraction.

\section{Comparison Table}

\begin{longtable}{p{3cm}p{3cm}p{3cm}p{3cm}p{3cm}}
	\toprule
	\textbf{Framework} & \textbf{Use Case}     & \textbf{Quantisation}   & \textbf{Batching} & \textbf{GPU req.} \\
	\midrule
	HF Transformers    & Research, prototyping & GPTQ, AWQ, bitsandbytes & Static            & $\geq 1$ GPU      \\
	vLLM               & Production serving    & GPTQ, AWQ, FP8          & Continuous        & $\geq 1$ GPU      \\
	TensorRT-LLM       & Max throughput        & FP8, INT4, INT8         & In-flight         & $\geq 1$ GPU      \\
	llama.cpp          & Local/edge inference  & GGUF 2--8 bit           & Static            & CPU/any GPU       \\
	LangChain          & Application dev       & N/A (uses backend)      & N/A               & N/A               \\
	LlamaIndex         & RAG systems           & N/A (uses backend)      & N/A               & N/A               \\
	\bottomrule
\end{longtable}

\section{Exercises}

\begin{enumerate}
	\item Deploy a Llama 3.1 8B model using vLLM with tensor parallelism on 2 GPUs.
	\item Quantise a model to 4-bit using AWQ and compare the perplexity vs.\ FP16.
	\item Build a RAG pipeline with LlamaIndex that retrieves from a PDF collection.
	\item Compare the throughput of vLLM vs. llama.cpp for a 7B model on the same hardware.
\end{enumerate}

% ============================================================
\chapter{Evaluation and Advanced Topics}
\label{ch:evaluation}

\section{Benchmark Suites}

\begin{longtable}{p{3.5cm}p{5cm}p{5cm}}
	\toprule
	\textbf{Benchmark} & \textbf{What it measures}          & \textbf{Format}      \\
	\midrule
	MMLU               & Multi-task knowledge (57 subjects) & 4-way MCQ            \\
	HumanEval          & Code generation correctness        & Function completion  \\
	MT-Bench           & Multi-turn conversation quality    & LLM-as-judge scoring \\
	GSM8K              & Grade-school math reasoning        & Chain-of-thought     \\
	HELM               & Holistic evaluation (42 scenarios) & Multi-metric         \\
	BigBench           & Diverse reasoning tasks            & Varied               \\
	HellaSwag          & Commonsense reasoning              & Sentence completion  \\
	\bottomrule
\end{longtable}

\section{Hallucination and Calibration}

An LLM is \emph{calibrated} if its confidence matches its accuracy:
\begin{equation}
	\mathbb{P}(\text{correct} \mid p_{\text{max}} = c) = c \quad \forall c \in [0,1].
\end{equation}

Calibration is measured with expected calibration error (ECE):
\begin{equation}
	\text{ECE} = \sum_{m=1}^{M} \frac{|B_m|}{N} \left|\text{acc}(B_m) - \overline{\text{conf}}(B_m)\right|,
\end{equation}
where bins $B_m$ partition $[0,1]$.

\subsection{Hallucination Taxonomies}

\begin{itemize}
	\item \textbf{Factual hallucination:} model asserts false information (e.g., ``The Eiffel Tower is in London'').
	\item \textbf{Contextual hallucination:} model contradicts provided context (e.g., ignoring retrieved documents).
	\item \textbf{Reasoning hallucination:} model makes logical errors in multi-step reasoning.
	\item \textbf{Instruction hallucination:} model performs a different task than instructed.
\end{itemize}

Hallucination detection methods:
\begin{itemize}
	\item \textbf{Semantic entropy:} sample multiple answers, cluster by meaning, measure dispersion.
	\item \textbf{Self-consistency:} chain-of-thought samples reach the same answer.
	\item \textbf{P(True):} prompt the model to evaluate its own answer's correctness.
	\item \textbf{Context grounding:} verify that each claim is supported by the retrieved context.
\end{itemize}

\section{Safety Evaluation}

\begin{itemize}
	\item \textbf{Red-teaming:} adversarial testing to find failure modes.
	\item \textbf{Refusal rate:} fraction of harmful prompts correctly refused.
	\item \textbf{Sycophancy:} tendency to agree with user even when wrong.
	\item \textbf{Benchmark contamination:} test set overlap with training data inflates scores.
\end{itemize}

\section{Outlook: The Future of LLM Engineering}

The field is evolving rapidly. Key directions:
\begin{itemize}
	\item \textbf{MoE architecture:} Mixture-of-Experts enables larger models with constant inference cost.
	\item \textbf{Long context:} 1M+ token context windows are becoming practical.
	\item \textbf{Multimodality:} native image, video, and audio understanding.
	\item \textbf{Agentic workflows:} LLMs as autonomous agents with tool use and planning.
	\item \textbf{Hardware co-design:} custom accelerators (Groq, Cerebras, Etched) for transformer inference.
	\item \textbf{Test-time compute scaling:} spending more compute at inference (CoT, search, self-consistency) improves reasoning quality.
\end{itemize}

\section{Final Project}

\begin{enumerate}
	\item Select an open-source LLM and deploy it with vLLM.
	\item Implement a structured extraction pipeline that parses JSON outputs.
	\item Build an agent with tool calling (calculator, search, database).
	\item Measure latency, throughput, and TTFT under varying batch sizes.
	\item Quantise the model to 4-bit and measure the quality degradation.
	\item Write a report summarising the trade-offs observed.
\end{enumerate}

% ============================================================
\newpage
\begin{thebibliography}{99}

	\bibitem{shannon1948} C. E. Shannon, ``A Mathematical Theory of Communication,'' \emph{Bell System Technical Journal}, 1948.

	\bibitem{vaswani2017} A. Vaswani et al., ``Attention Is All You Need,'' \emph{NeurIPS}, 2017.

	\bibitem{kaplan2020} J. Kaplan et al., ``Scaling Laws for Neural Language Models,'' \emph{arXiv:2001.08361}, 2020.

	\bibitem{hoffmann2022} J. Hoffmann et al., ``Training Compute-Optimal Large Language Models,'' \emph{NeurIPS}, 2022.

	\bibitem{dao2022} T. Dao et al., ``FlashAttention: Fast and Memory-Efficient Exact Attention with IO-Awareness,'' \emph{NeurIPS}, 2022.

	\bibitem{leviathan2022} Y. Leviathan et al., ``Fast Inference from Transformers via Speculative Decoding,'' \emph{ICML}, 2022.

	\bibitem{chen2023} C. Chen et al., ``Accelerating Large Language Model Decoding with Speculative Sampling,'' \emph{arXiv:2302.01318}, 2023.

	\bibitem{frantar2023} E. Frantar et al., ``GPTQ: Accurate Post-Training Quantization for Generative Pre-Trained Transformers,'' \emph{ICLR}, 2023.

	\bibitem{lin2023} J. Lin et al., ``AWQ: Activation-aware Weight Quantization for LLM Compression and Acceleration,'' \emph{MLSys}, 2023.

	\bibitem{kwon2023} W. Kwon et al., ``Efficient Memory Management for Large Language Model Serving with PagedAttention,'' \emph{SOSP}, 2023.

	\bibitem{ouyang2022} L. Ouyang et al., ``Training Language Models to Follow Instructions with Human Feedback,'' \emph{NeurIPS}, 2022.

	\bibitem{rae2022} J. W. Rae et al., ``Scaling Language Models: Methods, Analysis \& Insights from Training Gopher,'' \emph{arXiv:2112.11446}, 2022.

	\bibitem{touvron2023} H. Touvron et al., ``Llama 2: Open Foundation and Fine-Tuned Chat Models,'' \emph{arXiv:2307.09288}, 2023.

	\bibitem{dubey2024} A. Dubey et al., ``The Llama 3 Herd of Models,'' \emph{arXiv:2407.21783}, 2024.

	\bibitem{jiang2023} A. Q. Jiang et al., ``Mistral 7B,'' \emph{arXiv:2310.06825}, 2023.

	\bibitem{tishby1999} N. Tishby et al., ``The Information Bottleneck Method,'' \emph{Allerton Conference}, 1999.

	\bibitem{perez2021} J. Perez et al., ``True Few-Shot Learning with Language Models,'' \emph{NeurIPS}, 2021.

	\bibitem{wei2022} J. Wei et al., ``Chain-of-Thought Prompting Elicits Reasoning in Large Language Models,'' \emph{NeurIPS}, 2022.

	\bibitem{lewis2020} P. Lewis et al., ``Retrieval-Augmented Generation for Knowledge-Intensive NLP Tasks,'' \emph{NeurIPS}, 2020.

	\bibitem{rafailov2023} R. Rafailov et al., ``Direct Preference Optimization: Your Language Model is Secretly a Reward Model,'' \emph{NeurIPS}, 2023.

	\bibitem{schulman2017} J. Schulman et al., ``Proximal Policy Optimization Algorithms,'' \emph{arXiv:1707.06347}, 2017.

	\bibitem{gu2023} A. Gu \& T. Dao, ``Mamba: Linear-Time Sequence Modeling with Selective State Spaces,'' \emph{arXiv:2312.00752}, 2023.

	\bibitem{dao2024} T. Dao \& A. Gu, ``Transformers are SSMs: Generalized Models and Efficient Algorithms,'' \emph{arXiv:2405.21060}, 2024.

	\bibitem{gu2022} A. Gu et al., ``Efficiently Modeling Long Sequences with Structured State Spaces,'' \emph{ICLR}, 2022.

	\bibitem{sennrich2016} R. Sennrich et al., ``Neural Machine Translation of Rare Words with Subword Units,'' \emph{ACL}, 2016.

	\bibitem{kudo2018} T. Kudo \& J. Richardson, ``SentencePiece: A Simple and Language Independent Subword Tokenizer,'' \emph{EMNLP}, 2018.

	\bibitem{chen2023pi} S. Chen et al., ``Extending Context Window of Large Language Models via Position Interpolation,'' \emph{arXiv:2306.15595}, 2023.

	\bibitem{peng2023} B. Peng et al., ``YaRN: Efficient Context Window Extension of Large Language Models,'' \emph{arXiv:2309.00071}, 2023.

	\bibitem{bloc2023} bloc97, ``NTK-Aware Scaled RoPE,'' \emph{GitHub issue \#133}, 2023.

	\bibitem{fedus2022} W. Fedus et al., ``Switch Transformers: Scaling to Trillion Parameter Models with Simple and Efficient Sparsity,'' \emph{JMLR}, 2022.

	\bibitem{lepikhin2021} D. Lepikhin et al., ``GShard: Scaling Giant Models with Conditional Computation and Automatic Sharding,'' \emph{ICLR}, 2021.

	\bibitem{dai2024} DeepSeek-AI, ``DeepSeek-V2: A Strong, Economical, and Efficient Mixture-of-Experts Language Model,'' \emph{arXiv:2405.04434}, 2024.

	\bibitem{dai2025} DeepSeek-AI, ``DeepSeek-R1: Incentivizing Reasoning Capability in LLMs via Reinforcement Learning,'' \emph{arXiv:2501.12948}, 2025.

	\bibitem{razin2024} N. Razin et al., ``Likelihood Displacement in Direct Preference Optimization,'' \emph{arXiv:2405.14853}, 2024.

	\bibitem{weiss2021} G. Weiss et al., ``Thinking Like a Transformer: A Theoretical Analysis of the Expressive Power of Transformers,'' \emph{ICLR}, 2021.

	\bibitem{merrill2023} W. Merrill \& A. Sabharwal, ``The Parallelism Tradeoff: Limitations of Log-Precision Transformers,'' \emph{ICLR}, 2023.

	\bibitem{press2022} O. Press et al., ``Train Short, Test Long: Attention with Linear Biases Enables Input Length Extrapolation,'' \emph{ICLR}, 2022.

	\bibitem{shazeer2017} N. Shazeer et al., ``Outrageously Large Neural Networks: The Sparsely-Gated Mixture-of-Experts Layer,'' \emph{ICLR}, 2017.

	\bibitem{zhou2024} Z. Zhou et al., ``LongRoPE: Extending LLM Context Window Beyond 2 Million Tokens,'' \emph{arXiv:2402.13753}, 2024.

	\bibitem{cai2024} T. Cai et al., ``Medusa: Simple LLM Inference Acceleration Framework with Multiple Decoding Heads,'' \emph{ICML}, 2024.

\end{thebibliography}

\end{document}
